\worldchapter{\trans{items_title}}{appendix-items}
\label{ch:appendix-items}

\begin{multicols}{2}


    \section{Erste Hilfe}



    \subsection*{Bandagen}

    Ermöglicht die Anwendung von Erster Hilfe.

    \begin{tabular}{@{}ll@{}}
        \textbf{\trans{weight_label}:} & 0.30 \\
        \textbf{\trans{price_label}:} & 5 \\
        \textbf{\trans{rarity_label}:} & Gewöhnlich \\
        \textbf{\trans{concealment_label}:} & 0 \\
    \end{tabular}

    \wbif{phasesix}{
        \vspace{3mm}
\faIcon{cube}\hspace{0.5em}    }{}


    \section{Tränke und Gifte}



    \subsection*{Trank der Macht}

    Nach der Einnahme spürt der Charakter, wie sich seine Muskeln verhärten und seine Sinne schärfen. Die Anzahl aller Würfel die der Spieler wirft werden verdoppelt. Der Trank hält 2W6 Minuten an. Der Anwender ignoriert zudem alle Mali durch Erschöpfung oder Wunden und erhält einen temporären Rüstungsschutz von 2 gegen normalen Schaden (2xR) physischen Schaden (dieser addiert sich zu vorhandener Rüstung).

Sobald die Wirkung nachlässt, fordert die Magie ihren Tribut: Der Charakter erleidet sofort 2W6+6 Runden Erschöpfung und ist während diesem Zeitraum benommen (+2 Malus auf alle Mindestwürfe).

    \begin{tabular}{@{}ll@{}}
        \textbf{\trans{weight_label}:} & 0.30 \\
        \textbf{\trans{price_label}:} & 1200 \\
        \textbf{\trans{rarity_label}:} & Selten \\
        \textbf{\trans{concealment_label}:} & 1 \\
    \end{tabular}

    \wbif{phasesix}{
        \vspace{3mm}
\faIcon{flask}\hspace{0.5em}\faIcon{magic}\hspace{0.5em}    }{}


    \subsection*{Einfacher Heiltrank}

    Bei der Anwendung werden 1W3 Wunden geheilt. Enthält 3 Anwendungen.

    \begin{tabular}{@{}ll@{}}
        \textbf{\trans{weight_label}:} & 0.20 \\
        \textbf{\trans{price_label}:} & 100 \\
        \textbf{\trans{rarity_label}:} & Außergewöhnlich \\
        \textbf{\trans{concealment_label}:} & 0 \\
        \textbf{\trans{charges_label}:} & 3\\
    \end{tabular}

    \wbif{phasesix}{
        \vspace{3mm}
\faIcon{magic}\hspace{0.5em}    }{}


    \subsection*{Phiole Flugschlangengift}

    Eine Phiole gefüllt mit Gift einer Flugschlange.

    \begin{tabular}{@{}ll@{}}
        \textbf{\trans{weight_label}:} & 0.20 \\
        \textbf{\trans{price_label}:} & 40 \\
        \textbf{\trans{rarity_label}:} & Selten \\
        \textbf{\trans{concealment_label}:} & 0 \\
    \end{tabular}

    \wbif{phasesix}{
        \vspace{3mm}
\faIcon{magic}\hspace{0.5em}    }{}


    \subsection*{Einfache Wundtinktur}

    Bei erfolgreicher Anwendung mit \textit{Erster Hilfe} und einer Bandage heilt die Bandage 1W3 Wunden zusätzlich.

    \begin{tabular}{@{}ll@{}}
        \textbf{\trans{weight_label}:} & 0.80 \\
        \textbf{\trans{price_label}:} & 30 \\
        \textbf{\trans{rarity_label}:} & Gewöhnlich \\
        \textbf{\trans{concealment_label}:} & 0 \\
    \end{tabular}

    \wbif{phasesix}{
        \vspace{3mm}
\faIcon{cube}\hspace{0.5em}    }{}


    \subsection*{Karaffe Zaubertrank}

    Stellt bei Anwendung 2 Arkana wieder her. Enthält 3 Anwendungen.

    \begin{tabular}{@{}ll@{}}
        \textbf{\trans{weight_label}:} & 1.00 \\
        \textbf{\trans{price_label}:} & 200 \\
        \textbf{\trans{rarity_label}:} & Selten \\
        \textbf{\trans{concealment_label}:} & 0 \\
        \textbf{\trans{charges_label}:} & 3\\
    \end{tabular}

    \wbif{phasesix}{
        \vspace{3mm}
\faIcon{magic}\hspace{0.5em}    }{}


    \subsection*{Lachtrank}


    \begin{tabular}{@{}ll@{}}
        \textbf{\trans{weight_label}:} & 1.00 \\
        \textbf{\trans{price_label}:} & 10 \\
        \textbf{\trans{rarity_label}:} & Selten \\
        \textbf{\trans{concealment_label}:} & 0 \\
    \end{tabular}

    \wbif{phasesix}{
        \vspace{3mm}
\faIcon{magic}\hspace{0.5em}    }{}


    \subsection*{Trank des Schutzes}

    Bei der Anwendung erhält der Charakter 1W3 Boost.

Enthält 3 Anwendungen.

    \begin{tabular}{@{}ll@{}}
        \textbf{\trans{weight_label}:} & 0.20 \\
        \textbf{\trans{price_label}:} & 80 \\
        \textbf{\trans{rarity_label}:} & Selten \\
        \textbf{\trans{concealment_label}:} & 0 \\
        \textbf{\trans{charges_label}:} & 3\\
    \end{tabular}

    \wbif{phasesix}{
        \vspace{3mm}
\faIcon{magic}\hspace{0.5em}    }{}


    \subsection*{Schlangenöl}

    Diese Tinktur wird gerne fälschlicherweise als Heiltrank verkauft. Bei der Anwendung stellt der Trank eine Wunde wieder her.

    \begin{tabular}{@{}ll@{}}
        \textbf{\trans{weight_label}:} & 0.30 \\
        \textbf{\trans{price_label}:} & 100 \\
        \textbf{\trans{rarity_label}:} & Gewöhnlich \\
        \textbf{\trans{concealment_label}:} & 0 \\
    \end{tabular}

    \wbif{phasesix}{
        \vspace{3mm}
\faIcon{umbrella}\hspace{0.5em}\faIcon{magic}\hspace{0.5em}    }{}


    \subsection*{Kinstarchel Sekret}

    Dieses Sekret wird aus den Knochen toter Kinstarchel gewonnen. Es vermag, wenn es mit einem Trank vermischt wird, diesen beim Wurf zur Explosion zu bringen. Eine Phiole oder Karaffe des geworfenen Tranks wirkt im Umkreis von 1W3 Metern, als würde der Trank eingenommen werden. Gleiches gilt für Gifte.

    \begin{tabular}{@{}ll@{}}
        \textbf{\trans{weight_label}:} & 0.20 \\
        \textbf{\trans{price_label}:} & 800 \\
        \textbf{\trans{rarity_label}:} & Selten \\
        \textbf{\trans{concealment_label}:} & 3 \\
    \end{tabular}

    \wbif{phasesix}{
        \vspace{3mm}
\faIcon{dragon}\hspace{0.5em}    }{}


    \subsection*{Elixier des saften Schlummers}

    Ein Schlaftrunk,der das anwendende Wesen in einen ruhigen, sanften Schlaf versetzt. 

Nach dem Konsum fällt der Charakter innerhalb von zehn Minuten in einen festen Schlaf, der mindestens fünf Stunden anhält. Wird der Charakter während dieser Zeit gestört oder versucht man, ihn zu wecken, muss er einen Wurf auf Resistenz ablegen. Bei einem Erfolg erwacht der Charakter augenblicklich. Bei einem Misserfolg bleibt der Schlaf erhalten, und der Wurf kann bei der nächsten Störung oder einem erneuten Weckversuch wiederholt werden. Wichtig ist jedoch: Wird der Charakter geschüttelt, angegriffen oder sehr lauten Geräuschen ausgesetzt, erwacht er sofort.

    \begin{tabular}{@{}ll@{}}
        \textbf{\trans{weight_label}:} & 0.13 \\
        \textbf{\trans{price_label}:} & 10 \\
        \textbf{\trans{rarity_label}:} & Gewöhnlich \\
        \textbf{\trans{concealment_label}:} & 0 \\
    \end{tabular}

    \wbif{phasesix}{
        \vspace{3mm}
\faIcon{flask}\hspace{0.5em}    }{}


    \subsection*{Trank der tiefen Ruhe}

    Der Trank versetzt den Nutzer in einen tiefen, erholenden Schlaf. Er gilt unter Magiekundigen als wahres Wundermittel, um die körperlichen Beeinträchtigungen einer durchzechten Nacht gänzlich zu revidieren.

Die Essenz führt dazu, dass der Charakter innerhalb von fünf Minuten in einen tiefen, zehnstündigen Schlaf fällt, wobei er W3 Wunden regeneriert. Um aus diesem tiefen Schlummer zu erwachen, muss der Charakter bei Störungen oder Weckversuchen zwei erfolgreiche Würfe auf Resistenz in Folge erzielen. Nur direkte körperliche Gewalt (Schütteln/Angriff) oder extrem laute Geräusche bewirken ein sofortiges Erwachen.

    \begin{tabular}{@{}ll@{}}
        \textbf{\trans{weight_label}:} & 0.30 \\
        \textbf{\trans{price_label}:} & 20 \\
        \textbf{\trans{rarity_label}:} & Außergewöhnlich \\
        \textbf{\trans{concealment_label}:} & 0 \\
    \end{tabular}

    \wbif{phasesix}{
        \vspace{3mm}
\faIcon{flask}\hspace{0.5em}    }{}


    \subsection*{Beturias ewige Ruhe}

    Dieses Gebräu zwingt den Körper in einen tiefen, 24-stündigen Scheintod, der selbst Angriffe ignoriert und wundersame Heilung schenkt.

Als ihre treuer Begleiter, der Bär Grumm, schwer verletzt war, kochte die zwergische Bardin Beturia das Elixier aus Monddornbeeren und Nachtschiefer-Staub. 

Nach der Einnahme fällt der Charakter sofort in einen 24-stündigen Scheintod-Zustand, wobei seine Körperfunktionen drastisch gesenkt werden und er 2W6 Wunden heilt. Um aus diesem beinahe komatösen Schlaf zu erwachen, benötigt der Charakter vier erfolgreiche Würfe auf Resistenz in Folge. Da Schütteln, laute Geräusche oder Angriffe ihn nicht wecken, kann er nur durch zwei erfolgreiche Würfe auf Untersuchen als lebend identifiziert werden; andernfalls wird er für tot gehalten.

    \begin{tabular}{@{}ll@{}}
        \textbf{\trans{weight_label}:} & 0.40 \\
        \textbf{\trans{price_label}:} & 200 \\
        \textbf{\trans{rarity_label}:} & Legendär \\
        \textbf{\trans{concealment_label}:} & 0 \\
    \end{tabular}

    \wbif{phasesix}{
        \vspace{3mm}
\faIcon{flask}\hspace{0.5em}    }{}


    \subsection*{Sud der seichten Ermächtigung}

    Der Spieler fügt seinem Würfelpool für die Dauer von 1W6 Minuten einen W6 hinzu. Dies gilt für alle Attributs-, Fertigkeits-, Kampf-, Magie-, Wissenswürfe, usw.
Nach Ablauf der Wirkungszeit ist der Spieler leicht reizbar und neigt zu Streitereien.

Eine bläulich, trübe, scharf riechende Brühe, die ursprünglich von den barbarischen Stämmen der nördlichen Steppen gebraut wurde, um sich vor Raubzügen gegen die südlichen Königreiche aufzuputschen.

    \begin{tabular}{@{}ll@{}}
        \textbf{\trans{weight_label}:} & 0.20 \\
        \textbf{\trans{price_label}:} & 100 \\
        \textbf{\trans{rarity_label}:} & Außergewöhnlich \\
        \textbf{\trans{concealment_label}:} & 0 \\
    \end{tabular}

    \wbif{phasesix}{
        \vspace{3mm}
\faIcon{flask}\hspace{0.5em}    }{}


    \subsection*{Elixier der elfischen Macht}

    Der Spieler fügt seinem Würfelpool für die Dauer von 2W6 Minuten  2W6 hinzu.
Nach Ablauf erleidet der Charakter den Zustand geschockt 2 und ist 2W6 Minuten Betäubt, da der Stoffwechsel nach dem Abklingen des hochpotenten Öls abrupt herunterfährt. Elfische Charaktere erleiden diese Nebenwirkung nicht.

Ein klares, leicht öliges Elixier, das in der Phiole perlmuttartig schimmert , was von dem Effekt des Schillerwal-Trans herrührt. Es riecht nicht nach Fisch, sondern überraschend frisch nach Ozeanbrise und süßen Blüten.

    \begin{tabular}{@{}ll@{}}
        \textbf{\trans{weight_label}:} & 0.30 \\
        \textbf{\trans{price_label}:} & 450 \\
        \textbf{\trans{rarity_label}:} & Selten \\
        \textbf{\trans{concealment_label}:} & 2 \\
    \end{tabular}

    \wbif{phasesix}{
        \vspace{3mm}
\faIcon{flask}\hspace{0.5em}    }{}


    \subsection*{grobschlächtige Spritze}

    gefüllt mit einer seltsamen Flüssigkeit

    \begin{tabular}{@{}ll@{}}
        \textbf{\trans{weight_label}:} & 1.00 \\
        \textbf{\trans{price_label}:} & 10 \\
        \textbf{\trans{rarity_label}:} & Einmalig \\
        \textbf{\trans{concealment_label}:} & 0 \\
    \end{tabular}

    \wbif{phasesix}{
        \vspace{3mm}
\faIcon{umbrella}\hspace{0.5em}\faIcon{industry}\hspace{0.5em}\faIcon{ankh}\hspace{0.5em}\faIcon{syringe}\hspace{0.5em}\faIcon{spider}\hspace{0.5em}    }{}


    \section{Werfbares}



    \subsection*{Wurfnetz}

    Das Wurfnetz kann im Kampf geworfen werden, um den Gegner im Netz zu fangen. 

Ist der Werfen-Wurf gelungen, so gilt der Gegner als gefangen. Er benötigt einen Geschick-Wurf, um sich aus dem Netz zu befreien (1 Aktion). Solange der Gegner im Netz gefangen ist, kann er sich nicht bewegen, alle Aktionen sind schwere Proben.

    \begin{tabular}{@{}ll@{}}
        \textbf{\trans{weight_label}:} & 1.00 \\
        \textbf{\trans{price_label}:} & 30 \\
        \textbf{\trans{rarity_label}:} & Gewöhnlich \\
        \textbf{\trans{concealment_label}:} & 0 \\
    \end{tabular}

    \wbif{phasesix}{
        \vspace{3mm}
\faIcon{cube}\hspace{0.5em}    }{}


    \subsection*{Silberner Ritualdolch}


    \begin{tabular}{@{}ll@{}}
        \textbf{\trans{weight_label}:} & 1.00 \\
        \textbf{\trans{price_label}:} & 10 \\
        \textbf{\trans{rarity_label}:} & Gewöhnlich \\
        \textbf{\trans{concealment_label}:} & 0 \\
    \end{tabular}

    \wbif{phasesix}{
        \vspace{3mm}
\faIcon{chess-rook}\hspace{0.5em}\faIcon{dragon}\hspace{0.5em}\faIcon{ankh}\hspace{0.5em}\faIcon{magic}\hspace{0.5em}    }{}


    \section{Gefäße}



    \subsection*{Keramikflasche}


    \begin{tabular}{@{}ll@{}}
        \textbf{\trans{weight_label}:} & 0.20 \\
        \textbf{\trans{price_label}:} & 10 \\
        \textbf{\trans{rarity_label}:} & Gewöhnlich \\
        \textbf{\trans{concealment_label}:} & 0 \\
    \end{tabular}

    \wbif{phasesix}{
        \vspace{3mm}
\faIcon{cube}\hspace{0.5em}    }{}


    \subsection*{Tinkturenbeutel}

    Ein Beutel, in der Regel aus Leinen, welcher um den Körper getragen werden kann. Im Inneren sind Fächer für Flaschen oder Tiegel abgetrennt.

    \begin{tabular}{@{}ll@{}}
        \textbf{\trans{weight_label}:} & 0.50 \\
        \textbf{\trans{price_label}:} & 15 \\
        \textbf{\trans{rarity_label}:} & Gewöhnlich \\
        \textbf{\trans{concealment_label}:} & 4 \\
    \end{tabular}

    \wbif{phasesix}{
        \vspace{3mm}
\faIcon{eye}\hspace{0.5em}\faIcon{chess-rook}\hspace{0.5em}\faIcon{umbrella}\hspace{0.5em}    }{}


    \subsection*{Phiole}

    Eine Phiole aus Glas

    \begin{tabular}{@{}ll@{}}
        \textbf{\trans{weight_label}:} & 0.10 \\
        \textbf{\trans{price_label}:} & 20 \\
        \textbf{\trans{rarity_label}:} & Gewöhnlich \\
        \textbf{\trans{concealment_label}:} & 0 \\
    \end{tabular}

    \wbif{phasesix}{
        \vspace{3mm}
\faIcon{cube}\hspace{0.5em}    }{}


    \subsection*{Ledertasche}


    \begin{tabular}{@{}ll@{}}
        \textbf{\trans{weight_label}:} & 0.80 \\
        \textbf{\trans{price_label}:} & 15 \\
        \textbf{\trans{rarity_label}:} & Gewöhnlich \\
        \textbf{\trans{concealment_label}:} & 2 \\
    \end{tabular}

    \wbif{phasesix}{
        \vspace{3mm}
\faIcon{cube}\hspace{0.5em}    }{}


    \subsection*{Tuchbeutel}

    Der Tuchbeutel kann genutzt werden um Gegenstände darin zu lagern oder zu transportieren.

    \begin{tabular}{@{}ll@{}}
        \textbf{\trans{weight_label}:} & 0.50 \\
        \textbf{\trans{price_label}:} & 5 \\
        \textbf{\trans{rarity_label}:} & Gewöhnlich \\
        \textbf{\trans{concealment_label}:} & 0 \\
    \end{tabular}

    \wbif{phasesix}{
        \vspace{3mm}
\faIcon{cube}\hspace{0.5em}    }{}


    \subsection*{Lederranzen}

    Ein bequem zu tragender Lederranzen, in dem Gegenstände verstaut werden können.

    \begin{tabular}{@{}ll@{}}
        \textbf{\trans{weight_label}:} & 2.00 \\
        \textbf{\trans{price_label}:} & 20 \\
        \textbf{\trans{rarity_label}:} & Gewöhnlich \\
        \textbf{\trans{concealment_label}:} & 1 \\
    \end{tabular}

    \wbif{phasesix}{
        \vspace{3mm}
\faIcon{cube}\hspace{0.5em}    }{}


    \subsection*{Pergamenthülle}

    Hierin sind deine Dokumente sicher! Eine lederne, wasserfeste Hülle um Pergamente oder Dokumente zu verstauen.

    \begin{tabular}{@{}ll@{}}
        \textbf{\trans{weight_label}:} & 0.20 \\
        \textbf{\trans{price_label}:} & 40 \\
        \textbf{\trans{rarity_label}:} & Gewöhnlich \\
        \textbf{\trans{concealment_label}:} & 0 \\
    \end{tabular}

    \wbif{phasesix}{
        \vspace{3mm}
\faIcon{cube}\hspace{0.5em}    }{}


    \subsection*{Sack}

    Ein Sack aus Leinen, groß genug um viele Gegenstände zu transportieren.

    \begin{tabular}{@{}ll@{}}
        \textbf{\trans{weight_label}:} & 1.00 \\
        \textbf{\trans{price_label}:} & 10 \\
        \textbf{\trans{rarity_label}:} & Gewöhnlich \\
        \textbf{\trans{concealment_label}:} & 1 \\
    \end{tabular}

    \wbif{phasesix}{
        \vspace{3mm}
\faIcon{cube}\hspace{0.5em}    }{}


    \subsection*{Korb}

    In diesem Korb kann man Gegenstände oder anderen Objekte transportieren.

    \begin{tabular}{@{}ll@{}}
        \textbf{\trans{weight_label}:} & 1.00 \\
        \textbf{\trans{price_label}:} & 10 \\
        \textbf{\trans{rarity_label}:} & Gewöhnlich \\
        \textbf{\trans{concealment_label}:} & 0 \\
    \end{tabular}

    \wbif{phasesix}{
        \vspace{3mm}
\faIcon{cube}\hspace{0.5em}    }{}


    \subsection*{Packsattel}

    Ein Packsattel für die Benutzung an einem Pferd.

    \begin{tabular}{@{}ll@{}}
        \textbf{\trans{weight_label}:} & 4.00 \\
        \textbf{\trans{price_label}:} & 30 \\
        \textbf{\trans{rarity_label}:} & Gewöhnlich \\
        \textbf{\trans{concealment_label}:} & 0 \\
    \end{tabular}

    \wbif{phasesix}{
        \vspace{3mm}
\faIcon{cube}\hspace{0.5em}    }{}


    \subsection*{Eimer}

    Ein 10l Eimer.

    \begin{tabular}{@{}ll@{}}
        \textbf{\trans{weight_label}:} & 0.70 \\
        \textbf{\trans{price_label}:} & 5 \\
        \textbf{\trans{rarity_label}:} & Gewöhnlich \\
        \textbf{\trans{concealment_label}:} & 0 \\
    \end{tabular}

    \wbif{phasesix}{
        \vspace{3mm}
\faIcon{cube}\hspace{0.5em}    }{}


    \subsection*{Glasflasche}

    Eine Glasflasche, die mit allem möglichen gefüllt werden kann.

    \begin{tabular}{@{}ll@{}}
        \textbf{\trans{weight_label}:} & 0.20 \\
        \textbf{\trans{price_label}:} & 5 \\
        \textbf{\trans{rarity_label}:} & Gewöhnlich \\
        \textbf{\trans{concealment_label}:} & 0 \\
    \end{tabular}

    \wbif{phasesix}{
        \vspace{3mm}
\faIcon{cube}\hspace{0.5em}    }{}


    \subsection*{Wasserfaß}

    Dieses Faß kann mit 20l Flüssigkeit gefüllt werden.

    \begin{tabular}{@{}ll@{}}
        \textbf{\trans{weight_label}:} & 5.00 \\
        \textbf{\trans{price_label}:} & 10 \\
        \textbf{\trans{rarity_label}:} & Gewöhnlich \\
        \textbf{\trans{concealment_label}:} & 1 \\
    \end{tabular}

    \wbif{phasesix}{
        \vspace{3mm}
\faIcon{cube}\hspace{0.5em}    }{}


    \subsection*{Jadeschatulle}


    \begin{tabular}{@{}ll@{}}
        \textbf{\trans{weight_label}:} & 0.50 \\
        \textbf{\trans{price_label}:} & 50 \\
        \textbf{\trans{rarity_label}:} & Außergewöhnlich \\
        \textbf{\trans{concealment_label}:} & 0 \\
    \end{tabular}

    \wbif{phasesix}{
        \vspace{3mm}
\faIcon{cube}\hspace{0.5em}    }{}


    \subsection*{Goldener Rubin besetzter Krug}


    \begin{tabular}{@{}ll@{}}
        \textbf{\trans{weight_label}:} & 1.00 \\
        \textbf{\trans{price_label}:} & 10 \\
        \textbf{\trans{rarity_label}:} & Gewöhnlich \\
        \textbf{\trans{concealment_label}:} & 0 \\
    \end{tabular}

    \wbif{phasesix}{
        \vspace{3mm}
\faIcon{chess-rook}\hspace{0.5em}\faIcon{dragon}\hspace{0.5em}\faIcon{ankh}\hspace{0.5em}\faIcon{magic}\hspace{0.5em}    }{}


    \subsection*{Jadeschatulle}

    Schirmt 50 Arkana Thanium ab

    \begin{tabular}{@{}ll@{}}
        \textbf{\trans{weight_label}:} & 1.00 \\
        \textbf{\trans{price_label}:} & 10 \\
        \textbf{\trans{rarity_label}:} & Selten \\
        \textbf{\trans{concealment_label}:} & 0 \\
        \textbf{\trans{charges_label}:} & 50\\
    \end{tabular}

    \wbif{phasesix}{
        \vspace{3mm}
\faIcon{chess-rook}\hspace{0.5em}\faIcon{dragon}\hspace{0.5em}\faIcon{ankh}\hspace{0.5em}\faIcon{magic}\hspace{0.5em}    }{}


    \subsection*{Köcher}

    Zum Halten von Pfeilen

    \begin{tabular}{@{}ll@{}}
        \textbf{\trans{weight_label}:} & 0.00 \\
        \textbf{\trans{price_label}:} & 0 \\
        \textbf{\trans{rarity_label}:} & Gewöhnlich \\
        \textbf{\trans{concealment_label}:} & 0 \\
    \end{tabular}

    \wbif{phasesix}{
        \vspace{3mm}
\faIcon{chess-rook}\hspace{0.5em}\faIcon{dragon}\hspace{0.5em}    }{}


    \section{Werkzeuge}



    \subsection*{Zirkel}

    Ein Zirkel kann für die Navigation oder geometrische Aufgaben benutzt werden.

    \begin{tabular}{@{}ll@{}}
        \textbf{\trans{weight_label}:} & 0.20 \\
        \textbf{\trans{price_label}:} & 30 \\
        \textbf{\trans{rarity_label}:} & Gewöhnlich \\
        \textbf{\trans{concealment_label}:} & 0 \\
    \end{tabular}

    \wbif{phasesix}{
        \vspace{3mm}
\faIcon{cube}\hspace{0.5em}    }{}


    \subsection*{Fallenwerkzeug}

    Werkzeug zum Anbringen oder Entschärfen von Fallen

    \begin{tabular}{@{}ll@{}}
        \textbf{\trans{weight_label}:} & 0.00 \\
        \textbf{\trans{price_label}:} & 0 \\
        \textbf{\trans{rarity_label}:} & Gewöhnlich \\
        \textbf{\trans{concealment_label}:} & 0 \\
    \end{tabular}

    \wbif{phasesix}{
        \vspace{3mm}
\faIcon{chess-rook}\hspace{0.5em}\faIcon{dragon}\hspace{0.5em}    }{}


    \subsection*{Mörser und Stößel}


    \begin{tabular}{@{}ll@{}}
        \textbf{\trans{weight_label}:} & 0.50 \\
        \textbf{\trans{price_label}:} & 5 \\
        \textbf{\trans{rarity_label}:} & Gewöhnlich \\
        \textbf{\trans{concealment_label}:} & 0 \\
    \end{tabular}

    \wbif{phasesix}{
        \vspace{3mm}
\faIcon{cube}\hspace{0.5em}    }{}


    \subsection*{Messer}

    Zum Verarbeiten von Tieren

    \begin{tabular}{@{}ll@{}}
        \textbf{\trans{weight_label}:} & 0.00 \\
        \textbf{\trans{price_label}:} & 0 \\
        \textbf{\trans{rarity_label}:} & Gewöhnlich \\
        \textbf{\trans{concealment_label}:} & 0 \\
    \end{tabular}

    \wbif{phasesix}{
        \vspace{3mm}
\faIcon{chess-rook}\hspace{0.5em}\faIcon{dragon}\hspace{0.5em}    }{}


    \subsection*{Reiselaboratorium}

    Eine robuste, abschließbare Holzkiste oder eine Art Koffer, gefüllt mit den wichtigsten, bruchsicheren Geräten: Ein kleiner Mörser, Probierlöffel aus Horn, verschließbare Lederbeutel für Pulver und ein paar kleine, mit Wachs versiegelbare Glaskolben.
Kann auf Reisen verwendet werden.

    \begin{tabular}{@{}ll@{}}
        \textbf{\trans{weight_label}:} & 6.00 \\
        \textbf{\trans{price_label}:} & 50 \\
        \textbf{\trans{rarity_label}:} & Selten \\
        \textbf{\trans{concealment_label}:} & 8 \\
    \end{tabular}

    \wbif{phasesix}{
        \vspace{3mm}
\faIcon{flask}\hspace{0.5em}    }{}


    \subsection*{Federkiel}

    Ein Federkiel zum Schreiben

    \begin{tabular}{@{}ll@{}}
        \textbf{\trans{weight_label}:} & 0.10 \\
        \textbf{\trans{price_label}:} & 15 \\
        \textbf{\trans{rarity_label}:} & Gewöhnlich \\
        \textbf{\trans{concealment_label}:} & 0 \\
    \end{tabular}

    \wbif{phasesix}{
        \vspace{3mm}
\faIcon{eye}\hspace{0.5em}\faIcon{chess-rook}\hspace{0.5em}\faIcon{umbrella}\hspace{0.5em}    }{}


    \subsection*{Senkblei}

    Ein Senkblei, um etwa die Tiefe von etwas abzuschätzen.

    \begin{tabular}{@{}ll@{}}
        \textbf{\trans{weight_label}:} & 0.30 \\
        \textbf{\trans{price_label}:} & 10 \\
        \textbf{\trans{rarity_label}:} & Gewöhnlich \\
        \textbf{\trans{concealment_label}:} & 0 \\
    \end{tabular}

    \wbif{phasesix}{
        \vspace{3mm}
\faIcon{cube}\hspace{0.5em}    }{}


    \subsection*{Pergament}

    Ein Blatt Pergament zum beschreiben

    \begin{tabular}{@{}ll@{}}
        \textbf{\trans{weight_label}:} & 0.01 \\
        \textbf{\trans{price_label}:} & 10 \\
        \textbf{\trans{rarity_label}:} & Gewöhnlich \\
        \textbf{\trans{concealment_label}:} & 0 \\
    \end{tabular}

    \wbif{phasesix}{
        \vspace{3mm}
\faIcon{eye}\hspace{0.5em}\faIcon{chess-rook}\hspace{0.5em}    }{}


    \subsection*{Abakus}

    Der Abakus ist eine einfache Rechenmaschine. Kommt er zur Anwendung sind alle Mechanik Würfe leichte Aürfe.

    \begin{tabular}{@{}ll@{}}
        \textbf{\trans{weight_label}:} & 0.70 \\
        \textbf{\trans{price_label}:} & 80 \\
        \textbf{\trans{rarity_label}:} & Gewöhnlich \\
        \textbf{\trans{concealment_label}:} & 0 \\
    \end{tabular}

    \wbif{phasesix}{
        \vspace{3mm}
\faIcon{cube}\hspace{0.5em}    }{}


    \subsection*{Kleiner Kessel}

    Ein kleiner Kessel aus Eisen

    \begin{tabular}{@{}ll@{}}
        \textbf{\trans{weight_label}:} & 1.00 \\
        \textbf{\trans{price_label}:} & 5 \\
        \textbf{\trans{rarity_label}:} & Gewöhnlich \\
        \textbf{\trans{concealment_label}:} & 0 \\
    \end{tabular}

    \wbif{phasesix}{
        \vspace{3mm}
\faIcon{cube}\hspace{0.5em}    }{}


    \subsection*{Siegelring}

    Ein Siegelring aus Gold

    \begin{tabular}{@{}ll@{}}
        \textbf{\trans{weight_label}:} & 0.00 \\
        \textbf{\trans{price_label}:} & 0 \\
        \textbf{\trans{rarity_label}:} & Gewöhnlich \\
        \textbf{\trans{concealment_label}:} & 0 \\
    \end{tabular}

    \wbif{phasesix}{
        \vspace{3mm}
\faIcon{chess-rook}\hspace{0.5em}\faIcon{dragon}\hspace{0.5em}    }{}


    \subsection*{Fesselseil}

    Dieses Fesselseil ist dafür geeignet feste Knoten zu knüpfen.

    \begin{tabular}{@{}ll@{}}
        \textbf{\trans{weight_label}:} & 1.00 \\
        \textbf{\trans{price_label}:} & 20 \\
        \textbf{\trans{rarity_label}:} & Gewöhnlich \\
        \textbf{\trans{concealment_label}:} & 0 \\
    \end{tabular}

    \wbif{phasesix}{
        \vspace{3mm}
\faIcon{cube}\hspace{0.5em}    }{}


    \subsection*{Kleiner Webrahmen}

    Ein kleiner Webrahmen, um auf der Reise gewebte Stoffe herstellen zu können.

    \begin{tabular}{@{}ll@{}}
        \textbf{\trans{weight_label}:} & 2.00 \\
        \textbf{\trans{price_label}:} & 20 \\
        \textbf{\trans{rarity_label}:} & Gewöhnlich \\
        \textbf{\trans{concealment_label}:} & 1 \\
    \end{tabular}

    \wbif{phasesix}{
        \vspace{3mm}
\faIcon{cube}\hspace{0.5em}    }{}


    \subsection*{Brechstange}

    Gordon Freeman knows how to use it

    \begin{tabular}{@{}ll@{}}
        \textbf{\trans{weight_label}:} & 1.00 \\
        \textbf{\trans{price_label}:} & 29 \\
        \textbf{\trans{rarity_label}:} & Gewöhnlich \\
        \textbf{\trans{concealment_label}:} & 1 \\
    \end{tabular}

    \wbif{phasesix}{
        \vspace{3mm}
\faIcon{cube}\hspace{0.5em}    }{}


    \subsection*{Schaufel}


    \begin{tabular}{@{}ll@{}}
        \textbf{\trans{weight_label}:} & 1.00 \\
        \textbf{\trans{price_label}:} & 30 \\
        \textbf{\trans{rarity_label}:} & Gewöhnlich \\
        \textbf{\trans{concealment_label}:} & 3 \\
    \end{tabular}

    \wbif{phasesix}{
        \vspace{3mm}
\faIcon{cube}\hspace{0.5em}    }{}


    \subsection*{Improvisierter Diedrich}


    \begin{tabular}{@{}ll@{}}
        \textbf{\trans{weight_label}:} & 0.01 \\
        \textbf{\trans{price_label}:} & 0 \\
        \textbf{\trans{rarity_label}:} & Gewöhnlich \\
        \textbf{\trans{concealment_label}:} & 0 \\
    \end{tabular}

    \wbif{phasesix}{
        \vspace{3mm}
\faIcon{cube}\hspace{0.5em}    }{}


    \subsection*{Nägel}

    Eine handvoll simpler Nägel

    \begin{tabular}{@{}ll@{}}
        \textbf{\trans{weight_label}:} & 0.05 \\
        \textbf{\trans{price_label}:} & 0 \\
        \textbf{\trans{rarity_label}:} & Gewöhnlich \\
        \textbf{\trans{concealment_label}:} & 0 \\
    \end{tabular}

    \wbif{phasesix}{
        \vspace{3mm}
\faIcon{cube}\hspace{0.5em}    }{}


    \subsection*{Kohlestifte}

    Kohlestifte können verwendet werden um auf Pergament oder Papier zu schreiben.

    \begin{tabular}{@{}ll@{}}
        \textbf{\trans{weight_label}:} & 0.30 \\
        \textbf{\trans{price_label}:} & 5 \\
        \textbf{\trans{rarity_label}:} & Gewöhnlich \\
        \textbf{\trans{concealment_label}:} & 0 \\
    \end{tabular}

    \wbif{phasesix}{
        \vspace{3mm}
\faIcon{cube}\hspace{0.5em}    }{}


    \subsection*{Obsidian Ritualdolch}


    \begin{tabular}{@{}ll@{}}
        \textbf{\trans{weight_label}:} & 1.00 \\
        \textbf{\trans{price_label}:} & 100 \\
        \textbf{\trans{rarity_label}:} & Außergewöhnlich \\
        \textbf{\trans{concealment_label}:} & 0 \\
    \end{tabular}

    \wbif{phasesix}{
        \vspace{3mm}
\faIcon{cube}\hspace{0.5em}    }{}


    \subsection*{Gerätschaften des Alchemisten}

    Ein Satz aus einem Alembik (Destillierkolben mit Helm und Auffanggefäß), feuerfeste Keramiktiegel, ein Satz Mörser und Stößel (aus Stein oder Messing), ein Blasebalg für die Flamme, sowie grundlegende Filtertücher und Pergament zum Notieren.
Muss stationär aufgebaut werden.

    \begin{tabular}{@{}ll@{}}
        \textbf{\trans{weight_label}:} & 10.00 \\
        \textbf{\trans{price_label}:} & 80 \\
        \textbf{\trans{rarity_label}:} & Außergewöhnlich \\
        \textbf{\trans{concealment_label}:} & 15 \\
    \end{tabular}

    \wbif{phasesix}{
        \vspace{3mm}
\faIcon{flask}\hspace{0.5em}    }{}


    \subsection*{Pfeife}

    Eine Pfeife zum Rauchen von Tabak und ähnlichem.

    \begin{tabular}{@{}ll@{}}
        \textbf{\trans{weight_label}:} & 0.10 \\
        \textbf{\trans{price_label}:} & 100 \\
        \textbf{\trans{rarity_label}:} & Gewöhnlich \\
        \textbf{\trans{concealment_label}:} & 0 \\
    \end{tabular}

    \wbif{phasesix}{
        \vspace{3mm}
\faIcon{cube}\hspace{0.5em}    }{}


    \subsection*{Kleine Pfanne}


    \begin{tabular}{@{}ll@{}}
        \textbf{\trans{weight_label}:} & 1.00 \\
        \textbf{\trans{price_label}:} & 5 \\
        \textbf{\trans{rarity_label}:} & Gewöhnlich \\
        \textbf{\trans{concealment_label}:} & 0 \\
    \end{tabular}

    \wbif{phasesix}{
        \vspace{3mm}
\faIcon{cube}\hspace{0.5em}    }{}


    \subsection*{Flaschenzug}

    Ein einfacher Flaschenzug. Ein Seil wird für den Betrieb benötigt. Der Flaschenzug kann 100kg heben.

    \begin{tabular}{@{}ll@{}}
        \textbf{\trans{weight_label}:} & 2.00 \\
        \textbf{\trans{price_label}:} & 40 \\
        \textbf{\trans{rarity_label}:} & Gewöhnlich \\
        \textbf{\trans{concealment_label}:} & 0 \\
    \end{tabular}

    \wbif{phasesix}{
        \vspace{3mm}
\faIcon{cube}\hspace{0.5em}    }{}


    \subsection*{Pinsel}

    Nutze diesen Pinsel um auf einer Leinwand zu malen.

    \begin{tabular}{@{}ll@{}}
        \textbf{\trans{weight_label}:} & 0.10 \\
        \textbf{\trans{price_label}:} & 5 \\
        \textbf{\trans{rarity_label}:} & Gewöhnlich \\
        \textbf{\trans{concealment_label}:} & 0 \\
    \end{tabular}

    \wbif{phasesix}{
        \vspace{3mm}
\faIcon{cube}\hspace{0.5em}    }{}


    \subsection*{Chiffrierbuch}

    Ein Buch mit Chiffren zur Verschlüsselung

    \begin{tabular}{@{}ll@{}}
        \textbf{\trans{weight_label}:} & 0.00 \\
        \textbf{\trans{price_label}:} & 0 \\
        \textbf{\trans{rarity_label}:} & Gewöhnlich \\
        \textbf{\trans{concealment_label}:} & 0 \\
    \end{tabular}

    \wbif{phasesix}{
        \vspace{3mm}
\faIcon{chess-rook}\hspace{0.5em}\faIcon{dragon}\hspace{0.5em}    }{}


    \subsection*{Dietriche}

    Wird ein Dietrich mit dem Wissen Schloss knacken verwendet, so wird ein leichter Wurf statt eines normalen Wurfes geworfen.

    \begin{tabular}{@{}ll@{}}
        \textbf{\trans{weight_label}:} & 0.20 \\
        \textbf{\trans{price_label}:} & 30 \\
        \textbf{\trans{rarity_label}:} & Gewöhnlich \\
        \textbf{\trans{concealment_label}:} & 0 \\
    \end{tabular}

    \wbif{phasesix}{
        \vspace{3mm}
\faIcon{cube}\hspace{0.5em}    }{}


    \subsection*{Tisch des Essenzenbrauers}

    Eine Erweiterung zur standard Alchemisten Ausrüstung, die kompliziertere Prozesse wie Trockendestillation (Retorte) oder präzises Erhitzen (Sandbad) erlaubt. Zeigt einen semi-professionellen Alchemisten an.
Muss stationär aufgebaut werden.

    \begin{tabular}{@{}ll@{}}
        \textbf{\trans{weight_label}:} & 20.00 \\
        \textbf{\trans{price_label}:} & 100 \\
        \textbf{\trans{rarity_label}:} & Außergewöhnlich \\
        \textbf{\trans{concealment_label}:} & 20 \\
    \end{tabular}

    \wbif{phasesix}{
        \vspace{3mm}
\faIcon{flask}\hspace{0.5em}    }{}


    \subsection*{Hammer}


    \begin{tabular}{@{}ll@{}}
        \textbf{\trans{weight_label}:} & 2.00 \\
        \textbf{\trans{price_label}:} & 30 \\
        \textbf{\trans{rarity_label}:} & Gewöhnlich \\
        \textbf{\trans{concealment_label}:} & 1 \\
    \end{tabular}

    \wbif{phasesix}{
        \vspace{3mm}
\faIcon{cube}\hspace{0.5em}    }{}


    \subsection*{Feuerstein und Zunder}


    \begin{tabular}{@{}ll@{}}
        \textbf{\trans{weight_label}:} & 0.30 \\
        \textbf{\trans{price_label}:} & 10 \\
        \textbf{\trans{rarity_label}:} & Gewöhnlich \\
        \textbf{\trans{concealment_label}:} & 0 \\
    \end{tabular}

    \wbif{phasesix}{
        \vspace{3mm}
\faIcon{chess-rook}\hspace{0.5em}    }{}


    \subsection*{Reisigbesen}

    Ein Besen. Man kann mit ihm fegen.

    \begin{tabular}{@{}ll@{}}
        \textbf{\trans{weight_label}:} & 2.00 \\
        \textbf{\trans{price_label}:} & 10 \\
        \textbf{\trans{rarity_label}:} & Gewöhnlich \\
        \textbf{\trans{concealment_label}:} & 1 \\
    \end{tabular}

    \wbif{phasesix}{
        \vspace{3mm}
\faIcon{cube}\hspace{0.5em}    }{}


    \subsection*{Stundenglas}

    Das Stundenglas kann benutzt werden um die Zeit genaug abzuschätzen.

    \begin{tabular}{@{}ll@{}}
        \textbf{\trans{weight_label}:} & 0.30 \\
        \textbf{\trans{price_label}:} & 50 \\
        \textbf{\trans{rarity_label}:} & Gewöhnlich \\
        \textbf{\trans{concealment_label}:} & 0 \\
    \end{tabular}

    \wbif{phasesix}{
        \vspace{3mm}
\faIcon{eye}\hspace{0.5em}\faIcon{chess-rook}\hspace{0.5em}\faIcon{umbrella}\hspace{0.5em}    }{}


    \subsection*{Schiefertafel}

    Auf dieser Schiefertafel kann man schreiben, und man kann das geschriebene jederzeit wegwischen.

    \begin{tabular}{@{}ll@{}}
        \textbf{\trans{weight_label}:} & 0.50 \\
        \textbf{\trans{price_label}:} & 10 \\
        \textbf{\trans{rarity_label}:} & Gewöhnlich \\
        \textbf{\trans{concealment_label}:} & 0 \\
    \end{tabular}

    \wbif{phasesix}{
        \vspace{3mm}
\faIcon{cube}\hspace{0.5em}    }{}


    \subsection*{Tintenfass}

    Ein sicher verschlossenes Tintenfass, welches Tinte für eine Feder oder einen Gänsekiel enthält.

    \begin{tabular}{@{}ll@{}}
        \textbf{\trans{weight_label}:} & 0.60 \\
        \textbf{\trans{price_label}:} & 10 \\
        \textbf{\trans{rarity_label}:} & Gewöhnlich \\
        \textbf{\trans{concealment_label}:} & 0 \\
        \textbf{\trans{charges_label}:} & 25\\
    \end{tabular}

    \wbif{phasesix}{
        \vspace{3mm}
\faIcon{cube}\hspace{0.5em}    }{}


    \section{Lichter}



    \subsection*{Fackel}


    \begin{tabular}{@{}ll@{}}
        \textbf{\trans{weight_label}:} & 0.20 \\
        \textbf{\trans{price_label}:} & 2 \\
        \textbf{\trans{rarity_label}:} & Gewöhnlich \\
        \textbf{\trans{concealment_label}:} & 0 \\
    \end{tabular}

    \wbif{phasesix}{
        \vspace{3mm}
\faIcon{cube}\hspace{0.5em}    }{}


    \subsection*{Laterne}


    \begin{tabular}{@{}ll@{}}
        \textbf{\trans{weight_label}:} & 1.00 \\
        \textbf{\trans{price_label}:} & 40 \\
        \textbf{\trans{rarity_label}:} & Gewöhnlich \\
        \textbf{\trans{concealment_label}:} & 1 \\
    \end{tabular}

    \wbif{phasesix}{
        \vspace{3mm}
\faIcon{cube}\hspace{0.5em}    }{}


    \subsection*{Kerze}

    Eine Kerze. Brennt etwa 8 Stunden.

    \begin{tabular}{@{}ll@{}}
        \textbf{\trans{weight_label}:} & 0.20 \\
        \textbf{\trans{price_label}:} & 5 \\
        \textbf{\trans{rarity_label}:} & Gewöhnlich \\
        \textbf{\trans{concealment_label}:} & 0 \\
    \end{tabular}

    \wbif{phasesix}{
        \vspace{3mm}
\faIcon{cube}\hspace{0.5em}    }{}


    \subsection*{Pechfackel}

    Die Pechfackel brennt für etwa 8 Stunden und erzeugt ein angenehmes, großflächiges Licht.

    \begin{tabular}{@{}ll@{}}
        \textbf{\trans{weight_label}:} & 0.50 \\
        \textbf{\trans{price_label}:} & 10 \\
        \textbf{\trans{rarity_label}:} & Gewöhnlich \\
        \textbf{\trans{concealment_label}:} & 0 \\
    \end{tabular}

    \wbif{phasesix}{
        \vspace{3mm}
\faIcon{cube}\hspace{0.5em}    }{}


    \subsection*{Öllampe}

    Die Öllampe verbreitet ein angenehmes, großflächiges Licht, und ist nicht so windanfällig wie eine Fackel.

    \begin{tabular}{@{}ll@{}}
        \textbf{\trans{weight_label}:} & 1.00 \\
        \textbf{\trans{price_label}:} & 30 \\
        \textbf{\trans{rarity_label}:} & Gewöhnlich \\
        \textbf{\trans{concealment_label}:} & 0 \\
    \end{tabular}

    \wbif{phasesix}{
        \vspace{3mm}
\faIcon{cube}\hspace{0.5em}    }{}


    \subsection*{Sturmlaterne}

    Die Sturmlaterne ist besonders resistent gegen Wind und Wetter. Sie verbreitet ein angenehmes Licht.

    \begin{tabular}{@{}ll@{}}
        \textbf{\trans{weight_label}:} & 1.00 \\
        \textbf{\trans{price_label}:} & 60 \\
        \textbf{\trans{rarity_label}:} & Gewöhnlich \\
        \textbf{\trans{concealment_label}:} & 0 \\
    \end{tabular}

    \wbif{phasesix}{
        \vspace{3mm}
\faIcon{cube}\hspace{0.5em}    }{}


    \section{Überwachung}



    \subsection*{Handschellen}


    \begin{tabular}{@{}ll@{}}
        \textbf{\trans{weight_label}:} & 0.50 \\
        \textbf{\trans{price_label}:} & 80 \\
        \textbf{\trans{rarity_label}:} & Gewöhnlich \\
        \textbf{\trans{concealment_label}:} & 0 \\
    \end{tabular}

    \wbif{phasesix}{
        \vspace{3mm}
\faIcon{eye}\hspace{0.5em}\faIcon{chess-rook}\hspace{0.5em}\faIcon{umbrella}\hspace{0.5em}    }{}


    \subsection*{Fernrohr}

    Alle \textit{Wahrnehmungswürfe}, die mit Hilfe des Fernrohrs gemacht werden sind einfache Proben.

    \begin{tabular}{@{}ll@{}}
        \textbf{\trans{weight_label}:} & 0.50 \\
        \textbf{\trans{price_label}:} & 80 \\
        \textbf{\trans{rarity_label}:} & Gewöhnlich \\
        \textbf{\trans{concealment_label}:} & 0 \\
    \end{tabular}

    \wbif{phasesix}{
        \vspace{3mm}
\faIcon{cube}\hspace{0.5em}    }{}


    \section{Trekking Ausrüstung}



    \subsection*{Angelhaken und Schnur}

    Eine simple Anlgerausrüstung.

    \begin{tabular}{@{}ll@{}}
        \textbf{\trans{weight_label}:} & 0.20 \\
        \textbf{\trans{price_label}:} & 10 \\
        \textbf{\trans{rarity_label}:} & Gewöhnlich \\
        \textbf{\trans{concealment_label}:} & 0 \\
    \end{tabular}

    \wbif{phasesix}{
        \vspace{3mm}
\faIcon{cube}\hspace{0.5em}    }{}


    \subsection*{Wasserschlauch}

    Ein 1 Liter Lederschlauch um Wasser zu transportieren.

    \begin{tabular}{@{}ll@{}}
        \textbf{\trans{weight_label}:} & 0.30 \\
        \textbf{\trans{price_label}:} & 20 \\
        \textbf{\trans{rarity_label}:} & Gewöhnlich \\
        \textbf{\trans{concealment_label}:} & 0 \\
    \end{tabular}

    \wbif{phasesix}{
        \vspace{3mm}
\faIcon{cube}\hspace{0.5em}    }{}


    \subsection*{Zelt}

    Ein großes 4-personen Zelt. Es ist ein wenig anstrengend aufzubauen, bietet aber Platz und Schutz für 4-5 Personen.

    \begin{tabular}{@{}ll@{}}
        \textbf{\trans{weight_label}:} & 5.00 \\
        \textbf{\trans{price_label}:} & 70 \\
        \textbf{\trans{rarity_label}:} & Gewöhnlich \\
        \textbf{\trans{concealment_label}:} & 1 \\
    \end{tabular}

    \wbif{phasesix}{
        \vspace{3mm}
\faIcon{cube}\hspace{0.5em}    }{}


    \subsection*{Hängematte}

    Diese Hängematte kann aufgespannt werden um eine bequeme Schlafstätte zu bieten.

    \begin{tabular}{@{}ll@{}}
        \textbf{\trans{weight_label}:} & 2.00 \\
        \textbf{\trans{price_label}:} & 20 \\
        \textbf{\trans{rarity_label}:} & Gewöhnlich \\
        \textbf{\trans{concealment_label}:} & 0 \\
    \end{tabular}

    \wbif{phasesix}{
        \vspace{3mm}
\faIcon{cube}\hspace{0.5em}    }{}


    \subsection*{Seil (3m)}


    \begin{tabular}{@{}ll@{}}
        \textbf{\trans{weight_label}:} & 3.00 \\
        \textbf{\trans{price_label}:} & 30 \\
        \textbf{\trans{rarity_label}:} & Gewöhnlich \\
        \textbf{\trans{concealment_label}:} & 2 \\
    \end{tabular}

    \wbif{phasesix}{
        \vspace{3mm}
\faIcon{cube}\hspace{0.5em}    }{}


    \subsection*{Wurfhaken}

    Ein Wurfhaken, dazu gedacht dorthin geworfen zu werden wo er sich festhaken kann. Idealerweise wird er zusammen mit einem daran festgebundenen Seil verwendet.

    \begin{tabular}{@{}ll@{}}
        \textbf{\trans{weight_label}:} & 2.00 \\
        \textbf{\trans{price_label}:} & 90 \\
        \textbf{\trans{rarity_label}:} & Gewöhnlich \\
        \textbf{\trans{concealment_label}:} & 1 \\
    \end{tabular}

    \wbif{phasesix}{
        \vspace{3mm}
\faIcon{cube}\hspace{0.5em}    }{}


    \subsection*{Rucksack}


    \begin{tabular}{@{}ll@{}}
        \textbf{\trans{weight_label}:} & 1.20 \\
        \textbf{\trans{price_label}:} & 100 \\
        \textbf{\trans{rarity_label}:} & Gewöhnlich \\
        \textbf{\trans{concealment_label}:} & 2 \\
    \end{tabular}

    \wbif{phasesix}{
        \vspace{3mm}
\faIcon{cube}\hspace{0.5em}    }{}


    \subsection*{Feuerstein und Stahl}

    Eine Möglichkeit ein Feuer zu entfachen. Ein wenig anstrengend, aber eine sehr sichere Methode.

    \begin{tabular}{@{}ll@{}}
        \textbf{\trans{weight_label}:} & 0.20 \\
        \textbf{\trans{price_label}:} & 5 \\
        \textbf{\trans{rarity_label}:} & Gewöhnlich \\
        \textbf{\trans{concealment_label}:} & 0 \\
    \end{tabular}

    \wbif{phasesix}{
        \vspace{3mm}
\faIcon{cube}\hspace{0.5em}    }{}


    \subsection*{Decke}


    \begin{tabular}{@{}ll@{}}
        \textbf{\trans{weight_label}:} & 1.00 \\
        \textbf{\trans{price_label}:} & 50 \\
        \textbf{\trans{rarity_label}:} & Gewöhnlich \\
        \textbf{\trans{concealment_label}:} & 1 \\
    \end{tabular}

    \wbif{phasesix}{
        \vspace{3mm}
\faIcon{cube}\hspace{0.5em}    }{}


    \subsection*{Gürteltaschen}

    Bequem zu erreichende Gürteltaschen. Etwa 4 davon können an einem Gürtel angebracht werden.

    \begin{tabular}{@{}ll@{}}
        \textbf{\trans{weight_label}:} & 0.30 \\
        \textbf{\trans{price_label}:} & 30 \\
        \textbf{\trans{rarity_label}:} & Gewöhnlich \\
        \textbf{\trans{concealment_label}:} & 0 \\
    \end{tabular}

    \wbif{phasesix}{
        \vspace{3mm}
\faIcon{cube}\hspace{0.5em}    }{}


    \subsection*{Brennglas}

    Eine Lupe, mit der unter anderem ein Feuer entzündet werden kann.

    \begin{tabular}{@{}ll@{}}
        \textbf{\trans{weight_label}:} & 0.20 \\
        \textbf{\trans{price_label}:} & 50 \\
        \textbf{\trans{rarity_label}:} & Gewöhnlich \\
        \textbf{\trans{concealment_label}:} & 0 \\
    \end{tabular}

    \wbif{phasesix}{
        \vspace{3mm}
\faIcon{cube}\hspace{0.5em}    }{}


    \subsection*{Zunderdose}

    Eine Zunderdose. Mit dem Inhalt kann ein Feuer bequem entzündet werden.

    \begin{tabular}{@{}ll@{}}
        \textbf{\trans{weight_label}:} & 0.10 \\
        \textbf{\trans{price_label}:} & 20 \\
        \textbf{\trans{rarity_label}:} & Gewöhnlich \\
        \textbf{\trans{concealment_label}:} & 0 \\
    \end{tabular}

    \wbif{phasesix}{
        \vspace{3mm}
\faIcon{eye}\hspace{0.5em}\faIcon{chess-rook}\hspace{0.5em}\faIcon{umbrella}\hspace{0.5em}    }{}


    \subsection*{Lampenöl}

    Ein Gefäß voller Lampenöl, um Sturmlaternen oder Öllampen nachzufüllen.

    \begin{tabular}{@{}ll@{}}
        \textbf{\trans{weight_label}:} & 1.00 \\
        \textbf{\trans{price_label}:} & 20 \\
        \textbf{\trans{rarity_label}:} & Gewöhnlich \\
        \textbf{\trans{concealment_label}:} & 0 \\
    \end{tabular}

    \wbif{phasesix}{
        \vspace{3mm}
\faIcon{cube}\hspace{0.5em}    }{}


    \subsection*{Strickleiter}

    Wenn die Strickleiter zusammengelegt ist, ist sie einfach zu verstauen. Ausgerollt bietet sie eine spontane Leiter über 8 Meter Höhe.

    \begin{tabular}{@{}ll@{}}
        \textbf{\trans{weight_label}:} & 2.00 \\
        \textbf{\trans{price_label}:} & 40 \\
        \textbf{\trans{rarity_label}:} & Gewöhnlich \\
        \textbf{\trans{concealment_label}:} & 0 \\
    \end{tabular}

    \wbif{phasesix}{
        \vspace{3mm}
\faIcon{cube}\hspace{0.5em}    }{}


    \subsection*{Lasso}

    Dieses Seil ist dafür gemacht ein Lasso zu knüpfen, um Tiere einzufangen.

    \begin{tabular}{@{}ll@{}}
        \textbf{\trans{weight_label}:} & 2.00 \\
        \textbf{\trans{price_label}:} & 20 \\
        \textbf{\trans{rarity_label}:} & Gewöhnlich \\
        \textbf{\trans{concealment_label}:} & 0 \\
    \end{tabular}

    \wbif{phasesix}{
        \vspace{3mm}
\faIcon{cube}\hspace{0.5em}    }{}


    \subsection*{Fischnetz}

    Mit diesem Netz kann man gut angeln.

    \begin{tabular}{@{}ll@{}}
        \textbf{\trans{weight_label}:} & 1.00 \\
        \textbf{\trans{price_label}:} & 10 \\
        \textbf{\trans{rarity_label}:} & Gewöhnlich \\
        \textbf{\trans{concealment_label}:} & 0 \\
    \end{tabular}

    \wbif{phasesix}{
        \vspace{3mm}
\faIcon{cube}\hspace{0.5em}    }{}


    \subsection*{Schlafsack}


    \begin{tabular}{@{}ll@{}}
        \textbf{\trans{weight_label}:} & 1.00 \\
        \textbf{\trans{price_label}:} & 50 \\
        \textbf{\trans{rarity_label}:} & Gewöhnlich \\
        \textbf{\trans{concealment_label}:} & 2 \\
    \end{tabular}

    \wbif{phasesix}{
        \vspace{3mm}
\faIcon{cube}\hspace{0.5em}    }{}


    \subsection*{Schneeschuhe}

    Dieses Paar Schneeschuhe kann benutzt werden um auf Schnee bequem und schnell zu laufen.

    \begin{tabular}{@{}ll@{}}
        \textbf{\trans{weight_label}:} & 1.00 \\
        \textbf{\trans{price_label}:} & 20 \\
        \textbf{\trans{rarity_label}:} & Gewöhnlich \\
        \textbf{\trans{concealment_label}:} & 0 \\
    \end{tabular}

    \wbif{phasesix}{
        \vspace{3mm}
\faIcon{cube}\hspace{0.5em}    }{}


    \subsection*{Trockenfleisch}

    Trockenfleisch ist durch Lufttrocknung haltbar gemachtes Fleisch, das aus rohem oder erhitztem Fleisch oder Fleischerzeugnissen hergestellt werden kann.

    \begin{tabular}{@{}ll@{}}
        \textbf{\trans{weight_label}:} & 0.10 \\
        \textbf{\trans{price_label}:} & 5 \\
        \textbf{\trans{rarity_label}:} & Gewöhnlich \\
        \textbf{\trans{concealment_label}:} & 0 \\
    \end{tabular}

    \wbif{phasesix}{
        \vspace{3mm}
\faIcon{cube}\hspace{0.5em}    }{}


    \subsection*{Kletterhaken}

    Ein Kletterhaken kann angebracht werden um Seile darin zu befestigen. Um ihn in den Fels zu schlagen kann ein Hammer verwendet werden.

    \begin{tabular}{@{}ll@{}}
        \textbf{\trans{weight_label}:} & 1.00 \\
        \textbf{\trans{price_label}:} & 5 \\
        \textbf{\trans{rarity_label}:} & Gewöhnlich \\
        \textbf{\trans{concealment_label}:} & 0 \\
    \end{tabular}

    \wbif{phasesix}{
        \vspace{3mm}
\faIcon{cube}\hspace{0.5em}    }{}


    \subsection*{Fell einer Hyäne}

    Das abgezogene Fell einer ausgewachsenen Hyäne.

    \begin{tabular}{@{}ll@{}}
        \textbf{\trans{weight_label}:} & 2.00 \\
        \textbf{\trans{price_label}:} & 20 \\
        \textbf{\trans{rarity_label}:} & Gewöhnlich \\
        \textbf{\trans{concealment_label}:} & 0 \\
    \end{tabular}

    \wbif{phasesix}{
        \vspace{3mm}
\faIcon{dragon}\hspace{0.5em}    }{}


    \section{Nahrung / Proviant}

    Nahrungsgegenstände zur versorgung der Hungrigen Helden und Heldinen


    \subsection*{Tabak}

    Bestes Langbodenblatt, grobschnitt, vollmundig.

    \begin{tabular}{@{}ll@{}}
        \textbf{\trans{weight_label}:} & 0.05 \\
        \textbf{\trans{price_label}:} & 15 \\
        \textbf{\trans{rarity_label}:} & Gewöhnlich \\
        \textbf{\trans{concealment_label}:} & 0 \\
        \textbf{\trans{charges_label}:} & 20\\
    \end{tabular}

    \wbif{phasesix}{
        \vspace{3mm}
\faIcon{cube}\hspace{0.5em}    }{}


    \subsection*{Edler Wein}

    Eine Flasche guten Weins.

    \begin{tabular}{@{}ll@{}}
        \textbf{\trans{weight_label}:} & 1.00 \\
        \textbf{\trans{price_label}:} & 80 \\
        \textbf{\trans{rarity_label}:} & Gewöhnlich \\
        \textbf{\trans{concealment_label}:} & 0 \\
        \textbf{\trans{charges_label}:} & 3\\
    \end{tabular}

    \wbif{phasesix}{
        \vspace{3mm}
\faIcon{cube}\hspace{0.5em}    }{}


    \subsection*{Bier}

    Kalt, kühl, lecker! Son frisches Bier, Junge, lecker. Kalt muss es sein, Junge!

    \begin{tabular}{@{}ll@{}}
        \textbf{\trans{weight_label}:} & 1.00 \\
        \textbf{\trans{price_label}:} & 1 \\
        \textbf{\trans{rarity_label}:} & Gewöhnlich \\
        \textbf{\trans{concealment_label}:} & 0 \\
    \end{tabular}

    \wbif{phasesix}{
        \vspace{3mm}
\faIcon{cube}\hspace{0.5em}    }{}


    \subsection*{Dörrfleisch}

    Trockenfleisch, nahrhaft und lange haltbar

    \begin{tabular}{@{}ll@{}}
        \textbf{\trans{weight_label}:} & 0.50 \\
        \textbf{\trans{price_label}:} & 1 \\
        \textbf{\trans{rarity_label}:} & Gewöhnlich \\
        \textbf{\trans{concealment_label}:} & 0 \\
        \textbf{\trans{charges_label}:} & 3\\
    \end{tabular}

    \wbif{phasesix}{
        \vspace{3mm}
\faIcon{cube}\hspace{0.5em}    }{}


    \subsection*{Eintopf}

    Eintopf aus verschiedenen Zutaten, alles, was der Koch gefunden hat. Er ist vielleicht etwas schwer zu transportieren, aber der Eintopf enthält sicher viele nahrhafte Zutaten.

    \begin{tabular}{@{}ll@{}}
        \textbf{\trans{weight_label}:} & 0.30 \\
        \textbf{\trans{price_label}:} & 5 \\
        \textbf{\trans{rarity_label}:} & Gewöhnlich \\
        \textbf{\trans{concealment_label}:} & 0 \\
    \end{tabular}

    \wbif{phasesix}{
        \vspace{3mm}
\faIcon{cube}\hspace{0.5em}    }{}


    \subsection*{Dörrobst}

    Verschiedenes gedörrtes Obst. Meist Rosinen, Feigen, Pflaumen, Äpfel, Birnen, Datteln oder Aprikosen.

    \begin{tabular}{@{}ll@{}}
        \textbf{\trans{weight_label}:} & 0.25 \\
        \textbf{\trans{price_label}:} & 0 \\
        \textbf{\trans{rarity_label}:} & Gewöhnlich \\
        \textbf{\trans{concealment_label}:} & 0 \\
    \end{tabular}

    \wbif{phasesix}{
        \vspace{3mm}
\faIcon{dragon}\hspace{0.5em}    }{}


    \subsection*{Brot}

    Ein Brot

    \begin{tabular}{@{}ll@{}}
        \textbf{\trans{weight_label}:} & 0.50 \\
        \textbf{\trans{price_label}:} & 10 \\
        \textbf{\trans{rarity_label}:} & Gewöhnlich \\
        \textbf{\trans{concealment_label}:} & 0 \\
    \end{tabular}

    \wbif{phasesix}{
        \vspace{3mm}
\faIcon{chess-rook}\hspace{0.5em}\faIcon{dragon}\hspace{0.5em}\faIcon{ankh}\hspace{0.5em}\faIcon{magic}\hspace{0.5em}    }{}


    \subsection*{Leib Brot}


    \begin{tabular}{@{}ll@{}}
        \textbf{\trans{weight_label}:} & 0.50 \\
        \textbf{\trans{price_label}:} & 1 \\
        \textbf{\trans{rarity_label}:} & Gewöhnlich \\
        \textbf{\trans{concealment_label}:} & 0 \\
        \textbf{\trans{charges_label}:} & 1\\
    \end{tabular}

    \wbif{phasesix}{
        \vspace{3mm}
\faIcon{chess-rook}\hspace{0.5em}\faIcon{dragon}\hspace{0.5em}\faIcon{ankh}\hspace{0.5em}\faIcon{magic}\hspace{0.5em}    }{}


    \subsection*{Rübenschnaps}

    Der Rübenschnaps ist ein typisches Erzeugnis der menschlichen Ackerländer in den gemäßigteren Zonen Tirakans, insbesondere in den westlichen und zentralen Königreichen der Menschen.

Er wird aus fermentierten und destillierten Zuckerrüben hergestellt.

Es handelt sich um einen billigen, aber hochprozentigen Fusel. Er dient in der Alchemie lediglich als Lösungsmittel und Hitzequelle, um die aggressiven oder flüchtigen Eigenschaften anderer Substanzen (wie die Galle der Flugechse) zu extrahieren. Aufgrund seiner Reinheit und seines Mangels an komplexen Inhaltsstoffen ist er ideal als einfache Basis für Massenträger.

Der Rübenschnaps ist meist klar oder leicht gelblich-trüb und riecht stechend nach Ethanol und einer unterschwelligen, erdigen Süße. Er brennt beim Trinken und hinterlässt einen starken, unangenehmen Nachgeschmack.

Der Rübenschnaps symbolisiert die Ausdauer und Pragmatik der Menschen in Tirakan. Während die Elfen ihr Lebendes Wasser und die Zwerge ihr Tiefen-Salz haben, verlassen sich die Menschen auf einfache, schnell verfügbare Lösungen.

Er ist ein Massenartikel und ein wichtiger Handelsgegenstand in den Grenzregionen und den Söldnerlagern. Ein Großteil des Preises einfacher Tränke entfällt auf die Destillation und den Transport, nicht auf die Zutat selbst.

Für die hochrangigen Alchemisten in den Akademien ist der Rübenschnaps ein Zeichen des Dilettantismus; sie bevorzugen raffiniertere, weniger aggressive Lösungsmittel. Die Dorfalchemisten hingegen nutzen ihn wegen seiner Effizienz und Verfügbarkeit.

    \begin{tabular}{@{}ll@{}}
        \textbf{\trans{weight_label}:} & 0.70 \\
        \textbf{\trans{price_label}:} & 6 \\
        \textbf{\trans{rarity_label}:} & Gewöhnlich \\
        \textbf{\trans{concealment_label}:} & 0 \\
    \end{tabular}

    \wbif{phasesix}{
        \vspace{3mm}
\faIcon{flask}\hspace{0.5em}    }{}


    \subsection*{Kichererbsen}

    Diese kleine, gelb-grüne Hülsenfrucht ist ein wahrlich wundervolles Nahrungsmittel.

Abgesehen von ihrer alchemistischen Verwendung als Basis für Lachtränke, ist sie eine hervorragende Speise für müde Wanderer.

Man findet sie oft auf den reichhaltigen Märkten in und um Al Bah Ji Ra. Ihre Beliebtheit verbreitet sich jedoch langsam in ganz Tirakan.

Der Verzehr einer Handvoll gerösteter Kicher-Erbsen vertreibt trübe Gedanken und lässt den Schatten im Gemüt für kurze Zeit schrumpfen. Vor allem wenn man sie mit Gewürzen wie Knoblauch, Kümmel oder Pfeffer würzt.

Auch pur sind sie reich an Proteinen, komplexen Ballaststoffen und dienen so als nahrhafte Ergänzung für jeden Speiseplan.

„Wenn der Eintopf zu bitter schmeckt und das Leben zu schwer wiegt, wirf eine Handvoll Kichererbsen hinein. Dein Magen wird es dir mit einem Glucksen danken.“
— Meisterkoch Alar al-Din des Zelts der Sieben Schleier von El Kurru

    \begin{tabular}{@{}ll@{}}
        \textbf{\trans{weight_label}:} & 50.00 \\
        \textbf{\trans{price_label}:} & 1 \\
        \textbf{\trans{rarity_label}:} & Gewöhnlich \\
        \textbf{\trans{concealment_label}:} & 0 \\
    \end{tabular}

    \wbif{phasesix}{
        \vspace{3mm}
\faIcon{flask}\hspace{0.5em}    }{}


    \subsection*{Honig}

    In den kultivierten Regionen wird der Honig von fleißigen Imkern gewonnen, die ihre Stöcke oft in der Nähe von Kräutergärten aufstellen.

Dieser Honig ist klar, beständig und hat einen feinen Geschmack nach Thymian oder Lavendel. Er kostet in der Regel etwa 2 Gulden pro 5 Oth und ist die verlässliche Basis für jeden Alchemisten, der die Bitterkeit von Kräuterextrakten mildern will.

Wer es wagemutiger mag, sucht in den Wäldern nach den Nestern der Wildbienen. Dieser Honig ist dunkler, zäher und oft mit Pollen oder kleinen Wachsstücken versetzt. Er hat eine kräftige, fast erdige Note. Man sagt, wilder Honig aus den Urwäldern besitze eine stärkere regenerative Kraft für die Stimme. Ideal für Barden, die nach einer langen Nacht in der Schenke ihre Kehle ölen müssen.

    \begin{tabular}{@{}ll@{}}
        \textbf{\trans{weight_label}:} & 5.00 \\
        \textbf{\trans{price_label}:} & 2 \\
        \textbf{\trans{rarity_label}:} & Gewöhnlich \\
        \textbf{\trans{concealment_label}:} & 0 \\
    \end{tabular}

    \wbif{phasesix}{
        \vspace{3mm}
\faIcon{flask}\hspace{0.5em}    }{}


    \section{Fahrzeuge}



    \subsection*{Einfacher Einspänner}

    Der einfache Einspänner ist ein kleines Gefährt, welches von einem Pferd gezogen wird.

    \begin{tabular}{@{}ll@{}}
        \textbf{\trans{weight_label}:} & 120.00 \\
        \textbf{\trans{price_label}:} & 400 \\
        \textbf{\trans{rarity_label}:} & Gewöhnlich \\
        \textbf{\trans{concealment_label}:} & 8 \\
    \end{tabular}

    \wbif{phasesix}{
        \vspace{3mm}
\faIcon{eye}\hspace{0.5em}\faIcon{chess-rook}\hspace{0.5em}\faIcon{umbrella}\hspace{0.5em}\faIcon{flag}\hspace{0.5em}    }{}


    \subsection*{Zweispännige Kutsche}

    Die Kutsche wird von zwei Pferden gezogen und hat optional eine Plane als Überdachung.

    \begin{tabular}{@{}ll@{}}
        \textbf{\trans{weight_label}:} & 220.00 \\
        \textbf{\trans{price_label}:} & 600 \\
        \textbf{\trans{rarity_label}:} & Gewöhnlich \\
        \textbf{\trans{concealment_label}:} & 8 \\
    \end{tabular}

    \wbif{phasesix}{
        \vspace{3mm}
\faIcon{eye}\hspace{0.5em}\faIcon{chess-rook}\hspace{0.5em}\faIcon{umbrella}\hspace{0.5em}\faIcon{flag}\hspace{0.5em}    }{}


    \subsection*{Vierspänner}

    Eine große, schwere Kutsche mit hölzernem Aufbau oder Plane. Sie wird von vier Pferden gezogen.

    \begin{tabular}{@{}ll@{}}
        \textbf{\trans{weight_label}:} & 400.00 \\
        \textbf{\trans{price_label}:} & 900 \\
        \textbf{\trans{rarity_label}:} & Gewöhnlich \\
        \textbf{\trans{concealment_label}:} & 10 \\
    \end{tabular}

    \wbif{phasesix}{
        \vspace{3mm}
\faIcon{eye}\hspace{0.5em}\faIcon{chess-rook}\hspace{0.5em}\faIcon{umbrella}\hspace{0.5em}\faIcon{flag}\hspace{0.5em}    }{}


    \subsection*{Rennkutsche}

    Die Rennkutsche ist besonders windschnittig.

    \begin{tabular}{@{}ll@{}}
        \textbf{\trans{weight_label}:} & 300.00 \\
        \textbf{\trans{price_label}:} & 1200 \\
        \textbf{\trans{rarity_label}:} & Gewöhnlich \\
        \textbf{\trans{concealment_label}:} & 10 \\
    \end{tabular}

    \wbif{phasesix}{
        \vspace{3mm}
\faIcon{eye}\hspace{0.5em}\faIcon{chess-rook}\hspace{0.5em}\faIcon{umbrella}\hspace{0.5em}\faIcon{flag}\hspace{0.5em}    }{}


    \subsection*{Streitwagen}

    Ein gut gearbeiteter Streitwagen bietet Schutz vor Angreifern und ermöglicht es, enge Wendungen zu fahren.

    \begin{tabular}{@{}ll@{}}
        \textbf{\trans{weight_label}:} & 500.00 \\
        \textbf{\trans{price_label}:} & 1000 \\
        \textbf{\trans{rarity_label}:} & Außergewöhnlich \\
        \textbf{\trans{concealment_label}:} & 10 \\
    \end{tabular}

    \wbif{phasesix}{
        \vspace{3mm}
\faIcon{eye}\hspace{0.5em}\faIcon{chess-rook}\hspace{0.5em}\faIcon{umbrella}\hspace{0.5em}\faIcon{flag}\hspace{0.5em}    }{}


    \subsection*{Hundeschlitten}

    Der Hundeschlitten wird von 8-10 Hunden gezogen und kann optional mit Reifen ausgestattet werden, um auf festem Untergrund zu fahren.

    \begin{tabular}{@{}ll@{}}
        \textbf{\trans{weight_label}:} & 80.00 \\
        \textbf{\trans{price_label}:} & 80 \\
        \textbf{\trans{rarity_label}:} & Gewöhnlich \\
        \textbf{\trans{concealment_label}:} & 0 \\
    \end{tabular}

    \wbif{phasesix}{
        \vspace{3mm}
\faIcon{eye}\hspace{0.5em}\faIcon{chess-rook}\hspace{0.5em}\faIcon{umbrella}\hspace{0.5em}\faIcon{flag}\hspace{0.5em}    }{}


    \subsection*{Ochsenkarren}

    Der Ochsenkarren wird von zwei Ochsen gezogen. Ein sehr langsames, aber zuverlässiges Transportmittel.

    \begin{tabular}{@{}ll@{}}
        \textbf{\trans{weight_label}:} & 250.00 \\
        \textbf{\trans{price_label}:} & 120 \\
        \textbf{\trans{rarity_label}:} & Gewöhnlich \\
        \textbf{\trans{concealment_label}:} & 10 \\
    \end{tabular}

    \wbif{phasesix}{
        \vspace{3mm}
\faIcon{eye}\hspace{0.5em}\faIcon{chess-rook}\hspace{0.5em}\faIcon{umbrella}\hspace{0.5em}\faIcon{flag}\hspace{0.5em}    }{}


    \subsection*{Planwagen}

    Ein von zwei Pferden gezogener Planwagen. Die Plane bietet Schutz vor den meisten Wettereinflüssen.

    \begin{tabular}{@{}ll@{}}
        \textbf{\trans{weight_label}:} & 400.00 \\
        \textbf{\trans{price_label}:} & 400 \\
        \textbf{\trans{rarity_label}:} & Gewöhnlich \\
        \textbf{\trans{concealment_label}:} & 10 \\
    \end{tabular}

    \wbif{phasesix}{
        \vspace{3mm}
\faIcon{eye}\hspace{0.5em}\faIcon{chess-rook}\hspace{0.5em}\faIcon{umbrella}\hspace{0.5em}\faIcon{flag}\hspace{0.5em}    }{}


    \subsection*{Kastenwagen}

    Der hölzerne Aufbau auf diesem Kastenwagen schützt vor Wind, Wetter und Einbrechern. Das Gefährt wird von einem Pferd gezogen.

    \begin{tabular}{@{}ll@{}}
        \textbf{\trans{weight_label}:} & 500.00 \\
        \textbf{\trans{price_label}:} & 600 \\
        \textbf{\trans{rarity_label}:} & Gewöhnlich \\
        \textbf{\trans{concealment_label}:} & 10 \\
    \end{tabular}

    \wbif{phasesix}{
        \vspace{3mm}
\faIcon{eye}\hspace{0.5em}\faIcon{chess-rook}\hspace{0.5em}\faIcon{umbrella}\hspace{0.5em}\faIcon{flag}\hspace{0.5em}    }{}


    \subsection*{Kanu}

    Das Kanu kann genutzt werden um Wasser zu überqueren. Es ist jedoch nicht hochseetauglich.

    \begin{tabular}{@{}ll@{}}
        \textbf{\trans{weight_label}:} & 20.00 \\
        \textbf{\trans{price_label}:} & 60 \\
        \textbf{\trans{rarity_label}:} & Gewöhnlich \\
        \textbf{\trans{concealment_label}:} & 8 \\
    \end{tabular}

    \wbif{phasesix}{
        \vspace{3mm}
\faIcon{cube}\hspace{0.5em}    }{}


    \subsection*{Kleines Ruderboot}

    Ein Ruderboot samt Rudern.

    \begin{tabular}{@{}ll@{}}
        \textbf{\trans{weight_label}:} & 100.00 \\
        \textbf{\trans{price_label}:} & 120 \\
        \textbf{\trans{rarity_label}:} & Gewöhnlich \\
        \textbf{\trans{concealment_label}:} & 8 \\
    \end{tabular}

    \wbif{phasesix}{
        \vspace{3mm}
\faIcon{cube}\hspace{0.5em}    }{}


    \section{Tierbedarf}



    \subsection*{Pferdefutter}

    Hochwertiges Pferdefutter, eine Portion reicht etwa für eine Woche

    \begin{tabular}{@{}ll@{}}
        \textbf{\trans{weight_label}:} & 1.00 \\
        \textbf{\trans{price_label}:} & 2 \\
        \textbf{\trans{rarity_label}:} & Gewöhnlich \\
        \textbf{\trans{concealment_label}:} & 0 \\
    \end{tabular}

    \wbif{phasesix}{
        \vspace{3mm}
\faIcon{cube}\hspace{0.5em}    }{}


    \subsection*{Tierfutter}

    Hochwertiges Tierfutter. Eine Portion reicht etwa eine Woche.

    \begin{tabular}{@{}ll@{}}
        \textbf{\trans{weight_label}:} & 1.00 \\
        \textbf{\trans{price_label}:} & 1 \\
        \textbf{\trans{rarity_label}:} & Gewöhnlich \\
        \textbf{\trans{concealment_label}:} & 0 \\
    \end{tabular}

    \wbif{phasesix}{
        \vspace{3mm}
\faIcon{cube}\hspace{0.5em}    }{}


    \subsection*{Zaumzeug}


    \begin{tabular}{@{}ll@{}}
        \textbf{\trans{weight_label}:} & 1.00 \\
        \textbf{\trans{price_label}:} & 70 \\
        \textbf{\trans{rarity_label}:} & Gewöhnlich \\
        \textbf{\trans{concealment_label}:} & 0 \\
    \end{tabular}

    \wbif{phasesix}{
        \vspace{3mm}
\faIcon{cube}\hspace{0.5em}    }{}


    \subsection*{Kummet}

    Ein gepolsterter Ring, welcher verwendet wird um Ochsen anzuschirren.

    \begin{tabular}{@{}ll@{}}
        \textbf{\trans{weight_label}:} & 1.00 \\
        \textbf{\trans{price_label}:} & 20 \\
        \textbf{\trans{rarity_label}:} & Gewöhnlich \\
        \textbf{\trans{concealment_label}:} & 0 \\
    \end{tabular}

    \wbif{phasesix}{
        \vspace{3mm}
\faIcon{cube}\hspace{0.5em}    }{}


    \subsection*{Pferdedecke}


    \begin{tabular}{@{}ll@{}}
        \textbf{\trans{weight_label}:} & 2.00 \\
        \textbf{\trans{price_label}:} & 40 \\
        \textbf{\trans{rarity_label}:} & Gewöhnlich \\
        \textbf{\trans{concealment_label}:} & 0 \\
    \end{tabular}

    \wbif{phasesix}{
        \vspace{3mm}
\faIcon{cube}\hspace{0.5em}    }{}


    \subsection*{Sattel}


    \begin{tabular}{@{}ll@{}}
        \textbf{\trans{weight_label}:} & 4.00 \\
        \textbf{\trans{price_label}:} & 80 \\
        \textbf{\trans{rarity_label}:} & Gewöhnlich \\
        \textbf{\trans{concealment_label}:} & 4 \\
    \end{tabular}

    \wbif{phasesix}{
        \vspace{3mm}
\faIcon{cube}\hspace{0.5em}    }{}


    \subsection*{Packsattel}

    Ein Sattel mit Taschen.

    \begin{tabular}{@{}ll@{}}
        \textbf{\trans{weight_label}:} & 5.00 \\
        \textbf{\trans{price_label}:} & 50 \\
        \textbf{\trans{rarity_label}:} & Gewöhnlich \\
        \textbf{\trans{concealment_label}:} & 4 \\
    \end{tabular}

    \wbif{phasesix}{
        \vspace{3mm}
\faIcon{cube}\hspace{0.5em}    }{}


    \subsection*{Striegelzeug}


    \begin{tabular}{@{}ll@{}}
        \textbf{\trans{weight_label}:} & 1.00 \\
        \textbf{\trans{price_label}:} & 30 \\
        \textbf{\trans{rarity_label}:} & Gewöhnlich \\
        \textbf{\trans{concealment_label}:} & 0 \\
    \end{tabular}

    \wbif{phasesix}{
        \vspace{3mm}
\faIcon{cube}\hspace{0.5em}    }{}


    \subsection*{Reitgerte}


    \begin{tabular}{@{}ll@{}}
        \textbf{\trans{weight_label}:} & 1.00 \\
        \textbf{\trans{price_label}:} & 20 \\
        \textbf{\trans{rarity_label}:} & Gewöhnlich \\
        \textbf{\trans{concealment_label}:} & 0 \\
    \end{tabular}

    \wbif{phasesix}{
        \vspace{3mm}
\faIcon{cube}\hspace{0.5em}    }{}


    \subsection*{Eisensporen}


    \begin{tabular}{@{}ll@{}}
        \textbf{\trans{weight_label}:} & 1.00 \\
        \textbf{\trans{price_label}:} & 10 \\
        \textbf{\trans{rarity_label}:} & Gewöhnlich \\
        \textbf{\trans{concealment_label}:} & 0 \\
    \end{tabular}

    \wbif{phasesix}{
        \vspace{3mm}
\faIcon{cube}\hspace{0.5em}    }{}


    \subsection*{Silbersporen}


    \begin{tabular}{@{}ll@{}}
        \textbf{\trans{weight_label}:} & 1.00 \\
        \textbf{\trans{price_label}:} & 50 \\
        \textbf{\trans{rarity_label}:} & Gewöhnlich \\
        \textbf{\trans{concealment_label}:} & 0 \\
    \end{tabular}

    \wbif{phasesix}{
        \vspace{3mm}
\faIcon{cube}\hspace{0.5em}    }{}


    \subsection*{Falknerhandschuh}


    \begin{tabular}{@{}ll@{}}
        \textbf{\trans{weight_label}:} & 2.00 \\
        \textbf{\trans{price_label}:} & 40 \\
        \textbf{\trans{rarity_label}:} & Gewöhnlich \\
        \textbf{\trans{concealment_label}:} & 0 \\
    \end{tabular}

    \wbif{phasesix}{
        \vspace{3mm}
\faIcon{cube}\hspace{0.5em}    }{}


    \subsection*{Maulkorb}


    \begin{tabular}{@{}ll@{}}
        \textbf{\trans{weight_label}:} & 1.00 \\
        \textbf{\trans{price_label}:} & 20 \\
        \textbf{\trans{rarity_label}:} & Gewöhnlich \\
        \textbf{\trans{concealment_label}:} & 0 \\
    \end{tabular}

    \wbif{phasesix}{
        \vspace{3mm}
\faIcon{cube}\hspace{0.5em}    }{}


    \subsection*{Halsband und Leine}

    Halsband und Leine für einen Hund. Oder den Lebensgefährten.

    \begin{tabular}{@{}ll@{}}
        \textbf{\trans{weight_label}:} & 1.00 \\
        \textbf{\trans{price_label}:} & 30 \\
        \textbf{\trans{rarity_label}:} & Gewöhnlich \\
        \textbf{\trans{concealment_label}:} & 0 \\
    \end{tabular}

    \wbif{phasesix}{
        \vspace{3mm}
\faIcon{cube}\hspace{0.5em}    }{}


    \subsection*{Vogelkäfig}


    \begin{tabular}{@{}ll@{}}
        \textbf{\trans{weight_label}:} & 1.00 \\
        \textbf{\trans{price_label}:} & 30 \\
        \textbf{\trans{rarity_label}:} & Gewöhnlich \\
        \textbf{\trans{concealment_label}:} & 5 \\
    \end{tabular}

    \wbif{phasesix}{
        \vspace{3mm}
\faIcon{cube}\hspace{0.5em}    }{}


    \section{Kuriositäten}



    \subsection*{Buch eines Akademieschülers}

    Ein Buch eines Schülers aus der Akademie. Es sind wiederkehrend Alchemistische Zeichen für Blut unbeholfen hineingezeichnet worden.

    \begin{tabular}{@{}ll@{}}
        \textbf{\trans{weight_label}:} & 1.00 \\
        \textbf{\trans{price_label}:} & 10 \\
        \textbf{\trans{rarity_label}:} & Einmalig \\
        \textbf{\trans{concealment_label}:} & 0 \\
    \end{tabular}

    \wbif{phasesix}{
        \vspace{3mm}
\faIcon{dragon}\hspace{0.5em}\faIcon{magic}\hspace{0.5em}    }{}


    \subsection*{Raststein}

    Gibt 1x Täglich die Möglichkeit 3W6 die entweder als Bonus Würfel für die Rast oder zum Wiederherstellen von Arkana genutzt werden können (bei 5 als Erfolg)

    \begin{tabular}{@{}ll@{}}
        \textbf{\trans{weight_label}:} & 0.20 \\
        \textbf{\trans{price_label}:} & 500 \\
        \textbf{\trans{rarity_label}:} & Gewöhnlich \\
        \textbf{\trans{concealment_label}:} & 0 \\
    \end{tabular}

    \wbif{phasesix}{
        \vspace{3mm}
\faIcon{magic}\hspace{0.5em}    }{}


    \subsection*{Handspiegel}

    Ein einfacher, kleiner Handspiegel

    \begin{tabular}{@{}ll@{}}
        \textbf{\trans{weight_label}:} & 0.30 \\
        \textbf{\trans{price_label}:} & 15 \\
        \textbf{\trans{rarity_label}:} & Gewöhnlich \\
        \textbf{\trans{concealment_label}:} & 0 \\
    \end{tabular}

    \wbif{phasesix}{
        \vspace{3mm}
\faIcon{cube}\hspace{0.5em}    }{}


    \subsection*{Scheiben der Puppen}


    \begin{tabular}{@{}ll@{}}
        \textbf{\trans{weight_label}:} & 1.00 \\
        \textbf{\trans{price_label}:} & 0 \\
        \textbf{\trans{rarity_label}:} & Einmalig \\
        \textbf{\trans{concealment_label}:} & 0 \\
        \textbf{\trans{charges_label}:} & 9\\
    \end{tabular}

    \wbif{phasesix}{
        \vspace{3mm}
\faIcon{chess-rook}\hspace{0.5em}\faIcon{dragon}\hspace{0.5em}\faIcon{ankh}\hspace{0.5em}\faIcon{magic}\hspace{0.5em}    }{}


    \subsection*{Shoggothenzahn}

    Zahn einer Oma-Shoggothe

    \begin{tabular}{@{}ll@{}}
        \textbf{\trans{weight_label}:} & 1.00 \\
        \textbf{\trans{price_label}:} & 10 \\
        \textbf{\trans{rarity_label}:} & Gewöhnlich \\
        \textbf{\trans{concealment_label}:} & 0 \\
    \end{tabular}

    \wbif{phasesix}{
        \vspace{3mm}
\faIcon{umbrella}\hspace{0.5em}\faIcon{industry}\hspace{0.5em}\faIcon{ankh}\hspace{0.5em}\faIcon{syringe}\hspace{0.5em}\faIcon{spider}\hspace{0.5em}    }{}


    \subsection*{Haarnadel}

    Kann auch als einfacher Dietrich und Stichwerkzeug dienen.

    \begin{tabular}{@{}ll@{}}
        \textbf{\trans{weight_label}:} & 0.03 \\
        \textbf{\trans{price_label}:} & 19 \\
        \textbf{\trans{rarity_label}:} & Gewöhnlich \\
        \textbf{\trans{concealment_label}:} & 0 \\
    \end{tabular}

    \wbif{phasesix}{
        \vspace{3mm}
\faIcon{chess-rook}\hspace{0.5em}\faIcon{umbrella}\hspace{0.5em}\faIcon{flag}\hspace{0.5em}\faIcon{plane}\hspace{0.5em}\faIcon{microchip}\hspace{0.5em}    }{}


    \subsection*{Ring, Silber}

    Ein silberner Ring

    \begin{tabular}{@{}ll@{}}
        \textbf{\trans{weight_label}:} & 0.10 \\
        \textbf{\trans{price_label}:} & 10 \\
        \textbf{\trans{rarity_label}:} & Gewöhnlich \\
        \textbf{\trans{concealment_label}:} & 0 \\
    \end{tabular}

    \wbif{phasesix}{
        \vspace{3mm}
\faIcon{cube}\hspace{0.5em}    }{}


    \subsection*{Stoffpuppe}

    Eine einfache Stoffpuppe.

    \begin{tabular}{@{}ll@{}}
        \textbf{\trans{weight_label}:} & 0.30 \\
        \textbf{\trans{price_label}:} & 10 \\
        \textbf{\trans{rarity_label}:} & Gewöhnlich \\
        \textbf{\trans{concealment_label}:} & 0 \\
    \end{tabular}

    \wbif{phasesix}{
        \vspace{3mm}
\faIcon{cube}\hspace{0.5em}    }{}


    \subsection*{Spielzeugäffchen}


    \begin{tabular}{@{}ll@{}}
        \textbf{\trans{weight_label}:} & 1.00 \\
        \textbf{\trans{price_label}:} & 10 \\
        \textbf{\trans{rarity_label}:} & Außergewöhnlich \\
        \textbf{\trans{concealment_label}:} & 0 \\
    \end{tabular}

    \wbif{phasesix}{
        \vspace{3mm}
\faIcon{dragon}\hspace{0.5em}    }{}


    \subsection*{Heiliges Symbol von Ravenkind}

    Heiliges Symbol des Rabenkindes
Wundersamer Gegenstand, legendär (erfordert die Einweihung durch einen Kleriker oder Paladin mit guter Gesinnung)

Das Heilige Symbol von Ravenkind ist ein einzigartiges heiliges Symbol, das den gutherzigen Gläubigen von Barovia heilig ist. Es entstand vor der Gründung einer Kirche in Barovia. Der Legende nach wurde es von einem riesigen Raben – oder einem Engel in Form eines riesigen Raben – an einen Paladin namens Lugdana übergeben. Lugdana nutzte das heilige Symbol, um bis zu ihrem Tod Vampirnester auszurotten und zu zerstören. Die Hohepriester von Ravenloft behielten und trugen das heilige Symbol nach Lugdanas Tod.

Das heilige Symbol ist ein Platinamulett in Form einer Sonne, in dessen Mitte ein großer Kristall eingebettet ist. Das heilige Symbol hat 10 Ladungen für die folgenden Eigenschaften. Es erhält täglich im Morgengrauen 1W6 + 4 Ladungen zurück.

Halten Sie Vampire. Als Aktion können Sie 1 Ladung ausgeben und das heilige Symbol präsentieren, damit es mit heiliger Kraft aufflammt. Vampire und Vampire, die innerhalb von 30 Fuß um das heilige Symbol erscheinen, müssen einen Weisheitsrettungswurf mit DC 15 durchführen, wenn es aufflammt. Bei einem fehlgeschlagenen Rettungswurf ist das Ziel 1 Minute lang gelähmt. Es kann den Rettungswurf am Ende jedes seiner Züge wiederholen, um den Effekt auf sich selbst zu beenden.

Verwandle Untote. Wenn Sie über die Funktion „Untote verwandeln“ oder „Unheilig machen“ verfügen, können Sie 3 Ladungen verbrauchen, wenn Sie bei Verwendung dieser Funktion das heilige Symbol präsentieren. Wenn Sie dies tun, haben Untote einen Nachteil bei ihren Rettungswürfen gegen den Effekt.

Sonnenlicht. Als Aktion können Sie beim Präsentieren des heiligen Symbols 5 Ladungen aufwenden, um es in einem Radius von 30 Fuß helles Licht und für weitere 30 Fuß gedämpftes Licht ausstrahlen zu lassen. Das Licht ist Sonnenlicht und dauert 10 Minuten oder bis die Wirkung beendet ist (keine Aktion erforderlich).

    \begin{tabular}{@{}ll@{}}
        \textbf{\trans{weight_label}:} & 1.00 \\
        \textbf{\trans{price_label}:} & 1 \\
        \textbf{\trans{rarity_label}:} & Einmalig \\
        \textbf{\trans{concealment_label}:} & 0 \\
        \textbf{\trans{charges_label}:} & 10\\
    \end{tabular}

    \wbif{phasesix}{
        \vspace{3mm}
\faIcon{dragon}\hspace{0.5em}    }{}


    \subsection*{Toraner Bürgerausweis}

    Dieser Ausweis erklärt den Träger zum Bürger Torans, und eröffnet ihm alle Rechte und Pflichten des Toraner Bürgertums.

    \begin{tabular}{@{}ll@{}}
        \textbf{\trans{weight_label}:} & 1.00 \\
        \textbf{\trans{price_label}:} & 0 \\
        \textbf{\trans{rarity_label}:} & Außergewöhnlich \\
        \textbf{\trans{concealment_label}:} & 0 \\
    \end{tabular}

    \wbif{phasesix}{
        \vspace{3mm}
\faIcon{dragon}\hspace{0.5em}    }{}


    \subsection*{Scherbe aus Tanium}


    \begin{tabular}{@{}ll@{}}
        \textbf{\trans{weight_label}:} & 1.00 \\
        \textbf{\trans{price_label}:} & 9999 \\
        \textbf{\trans{rarity_label}:} & Selten \\
        \textbf{\trans{concealment_label}:} & 0 \\
    \end{tabular}

    \wbif{phasesix}{
        \vspace{3mm}
\faIcon{dragon}\hspace{0.5em}    }{}


    \subsection*{Logbuch des Kapitäns}


    \begin{tabular}{@{}ll@{}}
        \textbf{\trans{weight_label}:} & 1.00 \\
        \textbf{\trans{price_label}:} & 10 \\
        \textbf{\trans{rarity_label}:} & Gewöhnlich \\
        \textbf{\trans{concealment_label}:} & 0 \\
    \end{tabular}

    \wbif{phasesix}{
        \vspace{3mm}
\faIcon{dice-six}\hspace{0.5em}\faIcon{umbrella}\hspace{0.5em}\faIcon{car}\hspace{0.5em}\faIcon{flask}\hspace{0.5em}\faIcon{spider}\hspace{0.5em}    }{}


    \subsection*{Haube des logischen Denkens}

    +1 Logik

    \begin{tabular}{@{}ll@{}}
        \textbf{\trans{weight_label}:} & 1.00 \\
        \textbf{\trans{price_label}:} & 1111 \\
        \textbf{\trans{rarity_label}:} & Legendär \\
        \textbf{\trans{concealment_label}:} & 0 \\
    \end{tabular}

    \wbif{phasesix}{
        \vspace{3mm}
\faIcon{cube}\hspace{0.5em}\faIcon{eye}\hspace{0.5em}\faIcon{dice-six}\hspace{0.5em}\faIcon{chess-rook}\hspace{0.5em}\faIcon{umbrella}\hspace{0.5em}\faIcon{flag}\hspace{0.5em}\faIcon{plane}\hspace{0.5em}\faIcon{microchip}\hspace{0.5em}\faIcon{space-shuttle}\hspace{0.5em}\faIcon{industry}\hspace{0.5em}\faIcon{syringe}\hspace{0.5em}\faIcon{9}\hspace{0.5em}\faIcon{book-open}\hspace{0.5em}\faIcon{dragon}\hspace{0.5em}\faIcon{ankh}\hspace{0.5em}\faIcon{magic}\hspace{0.5em}\faIcon{spider}\hspace{0.5em}\faIcon{fire}\hspace{0.5em}    }{}


    \subsection*{Ring, Gold}

    Ein goldener Ring.

    \begin{tabular}{@{}ll@{}}
        \textbf{\trans{weight_label}:} & 0.10 \\
        \textbf{\trans{price_label}:} & 60 \\
        \textbf{\trans{rarity_label}:} & Außergewöhnlich \\
        \textbf{\trans{concealment_label}:} & 0 \\
    \end{tabular}

    \wbif{phasesix}{
        \vspace{3mm}
\faIcon{cube}\hspace{0.5em}    }{}


    \subsection*{Talisman eines Sethlarn}

    Klaue eines Sethlarn an einem Lederband
Es geht eine enorme magische Macht von ihm aus.
Es stoppt den Alterungsprozess des Trägers.
Wird es abgelegt, holt die Zeit den Träger ein.

    \begin{tabular}{@{}ll@{}}
        \textbf{\trans{weight_label}:} & 0.30 \\
        \textbf{\trans{price_label}:} & 1000 \\
        \textbf{\trans{rarity_label}:} & Legendär \\
        \textbf{\trans{concealment_label}:} & 0 \\
    \end{tabular}

    \wbif{phasesix}{
        \vspace{3mm}
\faIcon{dragon}\hspace{0.5em}    }{}


    \subsection*{Sonnenuhr}

    Eine transportable Sonnenuhr.

    \begin{tabular}{@{}ll@{}}
        \textbf{\trans{weight_label}:} & 0.50 \\
        \textbf{\trans{price_label}:} & 20 \\
        \textbf{\trans{rarity_label}:} & Gewöhnlich \\
        \textbf{\trans{concealment_label}:} & 0 \\
    \end{tabular}

    \wbif{phasesix}{
        \vspace{3mm}
\faIcon{cube}\hspace{0.5em}    }{}


    \subsection*{Früchtekuchen}


    \begin{tabular}{@{}ll@{}}
        \textbf{\trans{weight_label}:} & 0.30 \\
        \textbf{\trans{price_label}:} & 10 \\
        \textbf{\trans{rarity_label}:} & Gewöhnlich \\
        \textbf{\trans{concealment_label}:} & 0 \\
    \end{tabular}

    \wbif{phasesix}{
        \vspace{3mm}
\faIcon{cube}\hspace{0.5em}    }{}


    \subsection*{Jonglierbälle}

    Entwerder du kannst es, oder du kannst es nicht.

    \begin{tabular}{@{}ll@{}}
        \textbf{\trans{weight_label}:} & 1.00 \\
        \textbf{\trans{price_label}:} & 10 \\
        \textbf{\trans{rarity_label}:} & Gewöhnlich \\
        \textbf{\trans{concealment_label}:} & 0 \\
    \end{tabular}

    \wbif{phasesix}{
        \vspace{3mm}
\faIcon{cube}\hspace{0.5em}    }{}


    \subsection*{Grimoire}

    Ein magisches Grimoire zur Aufzeichnung von Zaubern

    \begin{tabular}{@{}ll@{}}
        \textbf{\trans{weight_label}:} & 0.00 \\
        \textbf{\trans{price_label}:} & 0 \\
        \textbf{\trans{rarity_label}:} & Gewöhnlich \\
        \textbf{\trans{concealment_label}:} & 0 \\
    \end{tabular}

    \wbif{phasesix}{
        \vspace{3mm}
\faIcon{chess-rook}\hspace{0.5em}\faIcon{dragon}\hspace{0.5em}    }{}


    \subsection*{Brille}

    Eine Brille, hoffentlich auf deine Sehstärke abgestimmt.

    \begin{tabular}{@{}ll@{}}
        \textbf{\trans{weight_label}:} & 0.40 \\
        \textbf{\trans{price_label}:} & 80 \\
        \textbf{\trans{rarity_label}:} & Gewöhnlich \\
        \textbf{\trans{concealment_label}:} & 0 \\
    \end{tabular}

    \wbif{phasesix}{
        \vspace{3mm}
\faIcon{cube}\hspace{0.5em}    }{}


    \subsection*{Märchenbuch}

    Ein Buch mit Märchen.

    \begin{tabular}{@{}ll@{}}
        \textbf{\trans{weight_label}:} & 1.00 \\
        \textbf{\trans{price_label}:} & 10 \\
        \textbf{\trans{rarity_label}:} & Gewöhnlich \\
        \textbf{\trans{concealment_label}:} & 0 \\
    \end{tabular}

    \wbif{phasesix}{
        \vspace{3mm}
\faIcon{cube}\hspace{0.5em}    }{}


    \subsection*{Leuchtender Bovist}

    Ein grünlich leuchtender Pilz. Auch gepflückt glüht er weiter. Wird der Pilz erneut eingepflanzt, wächst Dieser weiter

    \begin{tabular}{@{}ll@{}}
        \textbf{\trans{weight_label}:} & 1.00 \\
        \textbf{\trans{price_label}:} & 10 \\
        \textbf{\trans{rarity_label}:} & Selten \\
        \textbf{\trans{concealment_label}:} & 0 \\
    \end{tabular}

    \wbif{phasesix}{
        \vspace{3mm}
\faIcon{chess-rook}\hspace{0.5em}\faIcon{dragon}\hspace{0.5em}\faIcon{ankh}\hspace{0.5em}\faIcon{flask}\hspace{0.5em}\faIcon{magic}\hspace{0.5em}    }{}


    \subsection*{Historische Bibel}

    Eine gebundene, historische Ausgabe der Bibel.

    \begin{tabular}{@{}ll@{}}
        \textbf{\trans{weight_label}:} & 1.00 \\
        \textbf{\trans{price_label}:} & 100 \\
        \textbf{\trans{rarity_label}:} & Gewöhnlich \\
        \textbf{\trans{concealment_label}:} & 0 \\
    \end{tabular}

    \wbif{phasesix}{
        \vspace{3mm}
\faIcon{chess-rook}\hspace{0.5em}\faIcon{umbrella}\hspace{0.5em}\faIcon{flag}\hspace{0.5em}\faIcon{plane}\hspace{0.5em}\faIcon{microchip}\hspace{0.5em}    }{}


    \subsection*{Goldenes Monokel}

    Ein goldenes Monokel, welches man zum Zweck der guten Sicht vor ein Auge einsetzen kann.

    \begin{tabular}{@{}ll@{}}
        \textbf{\trans{weight_label}:} & 1.00 \\
        \textbf{\trans{price_label}:} & 150 \\
        \textbf{\trans{rarity_label}:} & Gewöhnlich \\
        \textbf{\trans{concealment_label}:} & 0 \\
    \end{tabular}

    \wbif{phasesix}{
        \vspace{3mm}
\faIcon{umbrella}\hspace{0.5em}\faIcon{flag}\hspace{0.5em}\faIcon{dragon}\hspace{0.5em}\faIcon{fire}\hspace{0.5em}    }{}


    \subsection*{Die Flöte die alte Freunde ruft}


    \begin{tabular}{@{}ll@{}}
        \textbf{\trans{weight_label}:} & 1.00 \\
        \textbf{\trans{price_label}:} & 10 \\
        \textbf{\trans{rarity_label}:} & Einmalig \\
        \textbf{\trans{concealment_label}:} & 0 \\
    \end{tabular}

    \wbif{phasesix}{
        \vspace{3mm}
\faIcon{dragon}\hspace{0.5em}    }{}


    \section{Zutaten}



    \subsection*{Beifuß (Artemisia vulgaris)}

    Eine Beifuß Pflanze. Die Spitzen des Sprosses werden verwendet um den Gallenfluss und die Verdauung in Gang zu bringen.

    \begin{tabular}{@{}ll@{}}
        \textbf{\trans{weight_label}:} & 0.10 \\
        \textbf{\trans{price_label}:} & 5 \\
        \textbf{\trans{rarity_label}:} & Gewöhnlich \\
        \textbf{\trans{concealment_label}:} & 0 \\
    \end{tabular}

    \wbif{phasesix}{
        \vspace{3mm}
\faIcon{cube}\hspace{0.5em}\faIcon{flask}\hspace{0.5em}    }{}


    \subsection*{Minze}

    Die Minze ist in Tirakan weit verbreitet, doch unter jenen, die sich mit der Kräuterkunde befassen, gilt sie als unverzichtbares Grundnahrungsmittel für Geist und Körper.

Vorkommen \& Wuchs:
Man findet die Minze meist an feuchten, halbschattigen Orten. Sie wuchert an den Ufern von Bachläufen, in verwunschenen Waldlichtungen oder in den Kräutergärten weiser Heilerinnen. 

Sie gilt als magenberuhigend, wirkt kühlend und lindert Halsschmerzen. Sie findet verwendung in Salben, Tinkturen und Tees. Wird auch oft genutzt, um den strengen Geschmack von Wildfleisch zu mildern oder um billigem Dünnbier eine frische Note zu verleihen.

    \begin{tabular}{@{}ll@{}}
        \textbf{\trans{weight_label}:} & 0.10 \\
        \textbf{\trans{price_label}:} & 1 \\
        \textbf{\trans{rarity_label}:} & Gewöhnlich \\
        \textbf{\trans{concealment_label}:} & 0 \\
    \end{tabular}

    \wbif{phasesix}{
        \vspace{3mm}
\faIcon{flask}\hspace{0.5em}    }{}


    \subsection*{Eibisch (Althaea officinalis)}

    Von dieser Heilpflanze wird die Wurzel verwendet. Diese wird kalt angesetzt und muss ungefähr zwei Stunden ziehen. Erst nach dem Ziehen wird die Flüssigkeit gesiebt und anschließend erhitzt. Die Stoffe bieten Schutz für die Schleimhäute und wirken reizlindernd. Eine hilfreiche Heilpflanze bei Magen-Darm-Problemen und Husten.

    \begin{tabular}{@{}ll@{}}
        \textbf{\trans{weight_label}:} & 0.10 \\
        \textbf{\trans{price_label}:} & 3 \\
        \textbf{\trans{rarity_label}:} & Gewöhnlich \\
        \textbf{\trans{concealment_label}:} & 0 \\
    \end{tabular}

    \wbif{phasesix}{
        \vspace{3mm}
\faIcon{cube}\hspace{0.5em}\faIcon{flask}\hspace{0.5em}    }{}


    \subsection*{Kamille (Matricaria recutita)}

    Kamille zählt zu den ältesten Heilpflanzen und wurde bereits im Mittelalter angewandt. Die Blüten haben eine heilende und beruhigende Wirkung. Äußerlich kann Kamille bei Entzündungen des Zahnfleisches, der Haut oder der Schleimhaut angewendet werden. Innerlich eingenommen wirkt sie bei Magen-Darm-Erkrankungen. Spülen und die Inhalation sind ebenfalls weit verbreitet.

    \begin{tabular}{@{}ll@{}}
        \textbf{\trans{weight_label}:} & 0.10 \\
        \textbf{\trans{price_label}:} & 2 \\
        \textbf{\trans{rarity_label}:} & Gewöhnlich \\
        \textbf{\trans{concealment_label}:} & 0 \\
    \end{tabular}

    \wbif{phasesix}{
        \vspace{3mm}
\faIcon{cube}\hspace{0.5em}\faIcon{flask}\hspace{0.5em}    }{}


    \subsection*{Salbei (Salvia officinalis)}

    Die Blätter des Salbeis haben eine entzündungshemmende, schweißhemmende und zusammenziehende Wirkung. Ein Tee oder Spülungen sind bei Halsentzündungen oder auch Schweißausbrüchen empfehlenswert.

    \begin{tabular}{@{}ll@{}}
        \textbf{\trans{weight_label}:} & 0.10 \\
        \textbf{\trans{price_label}:} & 5 \\
        \textbf{\trans{rarity_label}:} & Gewöhnlich \\
        \textbf{\trans{concealment_label}:} & 0 \\
    \end{tabular}

    \wbif{phasesix}{
        \vspace{3mm}
\faIcon{cube}\hspace{0.5em}\faIcon{flask}\hspace{0.5em}    }{}


    \subsection*{Lachsunder Seetang (Alga asgorania)}

    Eine zähe, dunkelgrüne bis bräunliche Meeresalge mit breiten, ledrigen Blättern, die an den Küsten von Lachsund im kalten Meerwasser gedeiht.

Sie schmeckt extrem salzig und verströmt einen intensiven Duft nach frischer Gischt und Jod.

In der alchemistischen Naturkunde dient er als Stabilisator für den Magen. Seine Inhaltsstoffe beruhigen den Gleichgewichtssinn und binden Giftstoffe im Verdauungstrakt. Für den alltäglichen Verzehr ist er nicht geeignet

Das Gewächs wird in großen Mengen an den felsigen Küsten von Asgoran angespült oder gezielt von den maritimen Sammlern der Siederei Falke aus dem Meer gefischt, getrocknet und verarbeitet. Auf den Fischmärkten von Lachsund ist er körbeweise erhältlich.

    \begin{tabular}{@{}ll@{}}
        \textbf{\trans{weight_label}:} & 0.01 \\
        \textbf{\trans{price_label}:} & 2 \\
        \textbf{\trans{rarity_label}:} & Gewöhnlich \\
        \textbf{\trans{concealment_label}:} & 0 \\
    \end{tabular}

    \wbif{phasesix}{
        \vspace{3mm}
\faIcon{dragon}\hspace{0.5em}\faIcon{flask}\hspace{0.5em}    }{}


    \subsection*{Schmetterlingsdrachensekret}

    Geht man vorsichtig zu Werke, so lassen sich Schmetterlingsdrachen melken. Sie sondern ein gar seltsames Sekret ab, welches denjenigen, der es konsumiert sofort in einen Schlaf mit faszinierenden Träumen fallen lässt.

Wird der Trank verabreicht oder genommen, so schläft die konsumierende Person mindestens acht Stunden tief und fest. Für diese Zeit wird die doppelte Rast angewendet. Der Schlafende ist allenfalls durch echte Schmerzen zu wecken.

    \begin{tabular}{@{}ll@{}}
        \textbf{\trans{weight_label}:} & 0.10 \\
        \textbf{\trans{price_label}:} & 200 \\
        \textbf{\trans{rarity_label}:} & Selten \\
        \textbf{\trans{concealment_label}:} & 0 \\
    \end{tabular}

    \wbif{phasesix}{
        \vspace{3mm}
\faIcon{dragon}\hspace{0.5em}\faIcon{flask}\hspace{0.5em}\faIcon{spider}\hspace{0.5em}    }{}


    \subsection*{Panthomeischer Nadelzapfen}

    Panthomeische Nadelzapfen sind die verholzten, harzreichen Fruchtstände der seltenen Panthomeischen Gebirgskiefer.

Diese Zapfen sind von einer harten, fast metallisch wirkenden Schuppenschicht überzogen, die im Inneren einen tiefvioletten, hochgradig klebrigen Balsam einschließt. Sie besitzen keine inhärente magische Energie, sind jedoch in der Alchemie für ihre stark konservierenden, strukturfestigenden und wärmespenden Eigenschaften bekannt.

Diese Zapfen wachsen ausschließlich in den eisigen, sturmdurchpeitschten Höhen der Panthomeischen Gipfel, weit oberhalb der Baumgrenze im unwegsamen Hochland.

Die Ernte ist beschwerlich, da die Zapfen nur von den Jägern oder den Hil-JabalJarh' geschlagen werden müssen.

Gelegentlich bringen zwergische Gebirgskarawanen vereinzelte Exemplare auf die Märkte in ganz Tirakan.

    \begin{tabular}{@{}ll@{}}
        \textbf{\trans{weight_label}:} & 0.15 \\
        \textbf{\trans{price_label}:} & 35 \\
        \textbf{\trans{rarity_label}:} & Selten \\
        \textbf{\trans{concealment_label}:} & 0 \\
    \end{tabular}

    \wbif{phasesix}{
        \vspace{3mm}
\faIcon{dragon}\hspace{0.5em}\faIcon{flask}\hspace{0.5em}    }{}


    \subsection*{Felsenmoos}

    Das Tirakan-Moos, von den Bergvölkern oft „Schattenflor“ oder „Schattensamt“ genannt, wächst meist in der Nähe von Bergwyvern Kollonien. Es ist in Felsspalten und kleinen höhlen zu finden. Man erzählt sich, dass es möglicherweise die Luft in engen Höhlen reinigt.

Das Moos wächst in dichten, schwammartigen Polstern. Seine Farbe ist ein tiefes, fast unnatürliches Dunkelviolett, das bei Berührung oder Luftzug in einem sanften, pulsierenden Indigo leuchtet (Biolumineszenz).
Bei Berührung hinterlässt es einen leicht klebrigen, metallisch riechenden Film auf der Haut.

Seine Hauptfunktion in Tränken ist die Erdung. Es verhindert, dass Energien oder rohe Kräfte sich verflüchtigen.

Roh zerkaut wirkt es stark schmerzlindernd und fiebersenkend, führt jedoch bei Überdosierung zu einer gefährlichen Verlangsamung des Herzschlags.

Die Ältesten behaupten, das Moos würde die Flüstern der Berge aufsaugen. Wer sein Ohr lange genug an ein Moospolster presst, soll die Stimmen der Vorfahren hören können oder den Berg, der vor Hunger knurrt.

    \begin{tabular}{@{}ll@{}}
        \textbf{\trans{weight_label}:} & 0.10 \\
        \textbf{\trans{price_label}:} & 80 \\
        \textbf{\trans{rarity_label}:} & Selten \\
        \textbf{\trans{concealment_label}:} & 0 \\
    \end{tabular}

    \wbif{phasesix}{
        \vspace{3mm}
\faIcon{dragon}\hspace{0.5em}\faIcon{flask}\hspace{0.5em}    }{}


    \subsection*{Kräutermischung}

    Eine leckere Kräutermischung um Essen zu verfeinern.

    \begin{tabular}{@{}ll@{}}
        \textbf{\trans{weight_label}:} & 0.10 \\
        \textbf{\trans{price_label}:} & 5 \\
        \textbf{\trans{rarity_label}:} & Gewöhnlich \\
        \textbf{\trans{concealment_label}:} & 0 \\
        \textbf{\trans{charges_label}:} & 10\\
    \end{tabular}

    \wbif{phasesix}{
        \vspace{3mm}
\faIcon{cube}\hspace{0.5em}\faIcon{flask}\hspace{0.5em}    }{}


    \subsection*{Fett}

    Es ist Fett. Pflanzlichem oder tierischem Ursprung. Es dient zum braten, zum Verfeinern von Speisen oder zum Schmieren.

    \begin{tabular}{@{}ll@{}}
        \textbf{\trans{weight_label}:} & 1.00 \\
        \textbf{\trans{price_label}:} & 1 \\
        \textbf{\trans{rarity_label}:} & Gewöhnlich \\
        \textbf{\trans{concealment_label}:} & 0 \\
    \end{tabular}

    \wbif{phasesix}{
        \vspace{3mm}
\faIcon{flask}\hspace{0.5em}    }{}


    \subsection*{Kohle}

    Ein rein tiefschwarzer Stoff, der durch das kontrollierte Verkohlen von hartem Holz unter Luftabschluss gewonnen wird. Kohle.

Sie dient alchemistisch als universelles Filtriermedium. In Rezepte eingebunden oder zu feinstem Staub zerstoßen, saugt sie Verunreinigungen, Trübungen und unerwünschte Bitterstoffe wie ein Schwamm aus siedenden Flüssigkeiten.  

Sie ist überall in Tirakan bei Köhlern, Hufschmieden und gewöhnlichen Krämern in rauen Mengen für wenige Deut oder Gulden auf den städtischen Märkten erhältlich.

    \begin{tabular}{@{}ll@{}}
        \textbf{\trans{weight_label}:} & 1.00 \\
        \textbf{\trans{price_label}:} & 2 \\
        \textbf{\trans{rarity_label}:} & Gewöhnlich \\
        \textbf{\trans{concealment_label}:} & 0 \\
    \end{tabular}

    \wbif{phasesix}{
        \vspace{3mm}
\faIcon{dragon}\hspace{0.5em}\faIcon{flask}\hspace{0.5em}    }{}


    \subsection*{Silberdistel (Carlina acaulis)}

    Eine tief am Boden wachsende, widerstandsfähige Pflanze mit einer silbrig glänzenden, stacheligen Blütenkrone, die selbst bei eisigen Temperaturen nicht verkümmerte.

In der Alchemie gilt sie als reinigender Katalysator, dessen Essenz Gifte im Blut bindet, Entzündungen hemmt und die natürlichen Abwehrkräfte des Körpers stärkt.  

Sie gedeiht vor allem auf den kargen, sonnengefluteten Bergwiesen und felsigen Hängen des Toraner Hochlands. Man findet sie jedoch auch in gemässigteren Gebieten in ganz Tirakan.

    \begin{tabular}{@{}ll@{}}
        \textbf{\trans{weight_label}:} & 0.05 \\
        \textbf{\trans{price_label}:} & 3 \\
        \textbf{\trans{rarity_label}:} & Gewöhnlich \\
        \textbf{\trans{concealment_label}:} & 0 \\
    \end{tabular}

    \wbif{phasesix}{
        \vspace{3mm}
\faIcon{dragon}\hspace{0.5em}\faIcon{flask}\hspace{0.5em}    }{}


    \subsection*{Bernstein}

    Ein glatter, ovaler Bernstein mit einem warmen, goldenen Schimmer. Seine polierte Oberfläche ist leicht transparent und reflektiert das Licht auf faszinierende Weise. Der handgroße Stein wirkt durch seine geschwungene Form wie ein natürlicher Talisman.

    \begin{tabular}{@{}ll@{}}
        \textbf{\trans{weight_label}:} & 0.10 \\
        \textbf{\trans{price_label}:} & 50 \\
        \textbf{\trans{rarity_label}:} & Außergewöhnlich \\
        \textbf{\trans{concealment_label}:} & 0 \\
    \end{tabular}

    \wbif{phasesix}{
        \vspace{3mm}
\faIcon{cube}\hspace{0.5em}    }{}


    \subsection*{Brennnessel (Urtica dioica)}

    Brennnesseln haben eine entwässernde und entzündungshemmende Wirkung. Ein Tee aus den Blättern der Brennnessel verschafft Linderung bei Rheuma und Gicht.

    \begin{tabular}{@{}ll@{}}
        \textbf{\trans{weight_label}:} & 0.10 \\
        \textbf{\trans{price_label}:} & 2 \\
        \textbf{\trans{rarity_label}:} & Gewöhnlich \\
        \textbf{\trans{concealment_label}:} & 0 \\
    \end{tabular}

    \wbif{phasesix}{
        \vspace{3mm}
\faIcon{cube}\hspace{0.5em}    }{}


    \subsection*{Obsidian (Staub)}

    Obsidianstaub ist ein mundaner, mineralischer alchemistischer Grundstoff, der aus vulkanischem Glas gewonnen wird.

Er besitzt keinerlei innewohnende magische Kraft, ist jedoch aufgrund seiner extremen Härte, seiner scharfen Kanten im mikroskopischen Bereich und seiner enormen Hitzebeständigkeit ein unschätzbarer mechanischer Katalysator.

In der Alchemie wird er genutzt, um Unreinheiten und Bitterstoffe aus kochenden Suden zu binden und feine Pflanzenessenzen zu zwingen, sich abzusetzten.

Obsidian entsteht dort, wo feurige Magma rasch an der frischen Luft oder durch Wasser abkühlt. Man findet das vulkanische Glas in den tiefen, erstarrten Lavaseen der Unterwelt oder an den rauen, vulkanischen Bruchstellen des Hochlands.
Er kann in schweren Mörsern zu einem mehlfeinen, grauen Staub zerstoßen werden.

Die unangefochtenen Hauptlieferanten für den besten Obsidianstaub in Tirakan sind die Zwerge der Xordai. Insbesondere die Kolonie Kor-Vamak, die direkt in einem erstarrten Lavasee errichtet wurde, hat sich auf den Abbau und die Verarbeitung dieses Materials spezialisiert.

Das dort ansässige Syndikat exportiert den Staub über ihre gut bewachten Handelskarawanen bis an die Oberfläche.

    \begin{tabular}{@{}ll@{}}
        \textbf{\trans{weight_label}:} & 0.01 \\
        \textbf{\trans{price_label}:} & 20 \\
        \textbf{\trans{rarity_label}:} & Außergewöhnlich \\
        \textbf{\trans{concealment_label}:} & 0 \\
    \end{tabular}

    \wbif{phasesix}{
        \vspace{3mm}
\faIcon{dragon}\hspace{0.5em}\faIcon{flask}\hspace{0.5em}    }{}


    \subsection*{Schlüsselblume (Primula veris)}

    Die Schlüsselblume war im als Fruchtbarkeits- und Schutzmittel bekannt. Heute hilft ein Schlüsselwurzeltee gegen Erkältungen. Salbei und Fenchel verstärken die Wirkung.

    \begin{tabular}{@{}ll@{}}
        \textbf{\trans{weight_label}:} & 0.10 \\
        \textbf{\trans{price_label}:} & 5 \\
        \textbf{\trans{rarity_label}:} & Gewöhnlich \\
        \textbf{\trans{concealment_label}:} & 0 \\
    \end{tabular}

    \wbif{phasesix}{
        \vspace{3mm}
\faIcon{cube}\hspace{0.5em}    }{}


    \subsection*{Frostflechte}

    Ein zähes Gewächs, das ausschließlich in den nördlichen Steppen und den kältesten Regionen der Gebirge Tirakans gedeiht.

Sie wächst bevorzugt an der Baumgrenze und auf kargen, windgepeitschten Felsvorsprüngen, wo die Temperaturen selbst im Sommer kaum über den Gefrierpunkt steigen. Sie ist bei zwergischen Prospektoren aus dem Norden gut bekannt.

Die Flechte selbst besitzt keine eigene Magie, aber sie hat eine extreme kältebindende Eigenschaft. Sie absorbiert die arktische Kälte ihrer Umgebung und hält sie für bis zu 3 Tage gespeichert, wenn die Umgebungstemperatur nicht über 30° steigt. Innerhalb dieses Zeitraums beträgt die Temperatur der Pflanze weiterhin die des letzten Standorts.

In der Alchemie dient sie daher als Katalysator der Stabilisierung, der starke, flüchtige Substanzen (wie Blut oder hochprozentigen Alkohol) in einen Zustand des Kälteschocks versetzt.

Die Frostflechte erscheint als ein unauffälliges, dichtes Geflecht in Tiefblau oder Weißgrau. Sie liegt wie ein verkrusteter Teppich auf den Steinen und ist bei oberflächlicher Betrachtung kaum von eisüberzogenen Felsen oder von gefrorenem Moos zu unterscheiden. Sie besitzt keine Blätter und keine Blüten. 

Sie ist oft ein Handelsartikel in den südlichen Königreichen der Menschen, da sie dort nicht wächst, aber für einfache Heil- und Stärkungstränke unerlässlich ist.

    \begin{tabular}{@{}ll@{}}
        \textbf{\trans{weight_label}:} & 0.10 \\
        \textbf{\trans{price_label}:} & 30 \\
        \textbf{\trans{rarity_label}:} & Selten \\
        \textbf{\trans{concealment_label}:} & 0 \\
    \end{tabular}

    \wbif{phasesix}{
        \vspace{3mm}
\faIcon{flask}\hspace{0.5em}    }{}


    \subsection*{Phiole Trollblut}

    Eine Phiole gefüllt mit Trollblut.

    \begin{tabular}{@{}ll@{}}
        \textbf{\trans{weight_label}:} & 0.10 \\
        \textbf{\trans{price_label}:} & 30 \\
        \textbf{\trans{rarity_label}:} & Selten \\
        \textbf{\trans{concealment_label}:} & 0 \\
    \end{tabular}

    \wbif{phasesix}{
        \vspace{3mm}
\faIcon{flask}\hspace{0.5em}\faIcon{magic}\hspace{0.5em}    }{}


    \subsection*{Schafgarbe (Achillea millefolium)}

    Schafgarbe wird wegen ihrer blutstillenden Wirkung eingesetzt. Die Blüten und die Blätter enthalten Gerb-, Bitter- und Mineralstoffe. Das ätherische Öl der Pflanze wirkt entzündungshemmend und krampflösend.

    \begin{tabular}{@{}ll@{}}
        \textbf{\trans{weight_label}:} & 0.10 \\
        \textbf{\trans{price_label}:} & 3 \\
        \textbf{\trans{rarity_label}:} & Gewöhnlich \\
        \textbf{\trans{concealment_label}:} & 0 \\
    \end{tabular}

    \wbif{phasesix}{
        \vspace{3mm}
\faIcon{cube}\hspace{0.5em}    }{}


    \subsection*{Monddornbeere}

    Die Monddornbeere ist ein Geschenk der tiefsten Nacht. Sie wächst als bodendeckender Strauch, dessen zarte Ranken und tiefgrüne Blätter von markanten, kurzen Dornen geschützt werden.

Man findet sie ausschließlich an Orten, die selten das Licht der Sonne erblicken, meist tief in uralten Wäldern oder in der Nähe von feuchten Grotten- und Höhleneingängen. Ihre wahre Pracht entfaltet sie nur unter dem Licht des Vollmonds, wenn ihre kleinen, beerenartigen Früchte in einem geheimnisvollen, matten Blau leuchten, fast als hätten sie das Licht der Himmelskugel selbst verschluckt.

Die Beere ist berüchtigt für ihre starke sedierende Wirkung. In kleiner Dosis wirkt sie beruhigend, doch konzentriert in einem Trank, kann ihre Essenz den Geist betäuben und den Körper in einen Zustand tiefer, traumloser Stille versetzen. Das Sammeln ist riskant, da ihre Dornen bei Berührung einen temporären Juckreiz auslösen können.

    \begin{tabular}{@{}ll@{}}
        \textbf{\trans{weight_label}:} & 0.10 \\
        \textbf{\trans{price_label}:} & 20 \\
        \textbf{\trans{rarity_label}:} & Außergewöhnlich \\
        \textbf{\trans{concealment_label}:} & 0 \\
    \end{tabular}

    \wbif{phasesix}{
        \vspace{3mm}
\faIcon{flask}\hspace{0.5em}    }{}


    \subsection*{Kieselstein}

    Ein kleiner Stein, benutzbar als Schleudermunition.

    \begin{tabular}{@{}ll@{}}
        \textbf{\trans{weight_label}:} & 0.10 \\
        \textbf{\trans{price_label}:} & 0 \\
        \textbf{\trans{rarity_label}:} & Gewöhnlich \\
        \textbf{\trans{concealment_label}:} & 0 \\
    \end{tabular}

    \wbif{phasesix}{
        \vspace{3mm}
\faIcon{eye}\hspace{0.5em}\faIcon{chess-rook}\hspace{0.5em}\faIcon{umbrella}\hspace{0.5em}\faIcon{flask}\hspace{0.5em}    }{}


    \subsection*{Gemeines Blutkraut}

    Das Blutkraut ist ein bodennahes Gewächs, das vor allem durch seine fleischigen, tiefroten Blätter auffällt.

Die feinen Adern auf der Blattoberfläche leuchten in einem kräftigen Scharlachrot, fast so, als würde echtes Blut durch sie pulsieren. Wenn man ein Blatt zerreibt, tritt ein klebriger, süßlich riechender Saft aus, der die Finger tagelang verfärbt. Es ist keine magische Pflanze im klassischen Sinne. Sie bezieht ihre Kraft aus dem eisenreichen Boden und dem fahlen Licht der dichten Wälder.

Ihr findet das Blutkraut meist dort, wo es schattig und feucht ist. Es bevorzugt den Fuß alter Eichen oder die unmittelbare Nähe von verrottendem Unterholz in tiefen Wäldern. Ein unaufmerksamer Reisender hält es oft für gewöhnlichen Purpur-Ampfer, doch ein geschulter Alchemist achtet auf die kleinen, perlenartigen Tautropfen, die sich stets an den Rändern der Blätter sammeln.

Man erntet es am besten in den frühen Morgenstunden, bevor die Sonne das Blätterdach durchbricht. Man schneidet nur die äußeren Blätter ab, um die Wurzel zu schonen.

    \begin{tabular}{@{}ll@{}}
        \textbf{\trans{weight_label}:} & 0.50 \\
        \textbf{\trans{price_label}:} & 5 \\
        \textbf{\trans{rarity_label}:} & Selten \\
        \textbf{\trans{concealment_label}:} & 0 \\
    \end{tabular}

    \wbif{phasesix}{
        \vspace{3mm}
\faIcon{flask}\hspace{0.5em}    }{}


    \subsection*{Zitronen-Ingwer (Zingiber meridionale)}

    Eine knollige, faserige Wurzel von hellbrauner Farbe, deren saftiges, gelbes Fleisch beim Aufschneiden eine brennende Schärfe und ein intensiv frisches Zitrusaroma freisetzt.

In der Alchemie gilt er als feuriger, die Durchblutung anregender Vitalisator. Er wird genutzt, um die Säfte des Körpers in Wallung zu bringen, Magenbeschwerden zu lindern und die Aufnahme von Heilelixieren im Blut drastisch zu beschleunigen.

Die Pflanze liebt feuchte Wärme und gedeiht prächtig in den fruchtbaren, flussreichen Ländern der O'grut oder wird über die weitreichenden Handelsnetzwerke der Menschen aus den südlichen Oasen nach Meridian importiert.

    \begin{tabular}{@{}ll@{}}
        \textbf{\trans{weight_label}:} & 0.50 \\
        \textbf{\trans{price_label}:} & 4 \\
        \textbf{\trans{rarity_label}:} & Gewöhnlich \\
        \textbf{\trans{concealment_label}:} & 0 \\
    \end{tabular}

    \wbif{phasesix}{
        \vspace{3mm}
\faIcon{dragon}\hspace{0.5em}\faIcon{flask}\hspace{0.5em}    }{}


    \subsection*{Schuppe einer Flussjungfer}

    Die Flussjungfer ist eine überaus scheue Kreatur, die, so munkelt man in Alchemistischen Kreisen, wohl noch kein sterbliches Wesen in Tirakan wahrhaftig zu Gesicht bekommen hat. Der einzige greifbare Beweis ihrer Existenz sind die Schuppen.

Man findet sie gelegentlich an schlammigen Flussufern, auf gischtumspülten Felsen oder in der Finsternis unterirdischer Seen entdeckt.

In der Dämmerung wirken diese Schuppen oft schlicht grün-bräunlich, fast wie herkömmliches Horn oder vertrocknetes Laub. Doch sobald die Sonne Tirakans die Oberfläche küsst, erwachen sie zum Leben und schimmern in allen Farben des Perlmutt. Sie erreichen oft die Größe einer stolzen Handfläche und ähneln in ihrer Beschaffenheit den Schuppen eines großen Fisches.

Glücklicherweise ist diese magische Zutat auf den Märkten des Reichs, besonders auf den Handelsrouten zwischen Toran und Yavon bis hinunter nach Meridian, kein seltener Anblick. Da sie regelmäßig gefunden werden, bleibt ihr Wert überschaubar, was sie zu einer ehrlichen Zutat macht. Sie sind sogar so verbreitet, dass es sich für Fälscher kaum lohnt, minderwertige Kopien zu fertigen. Aber ein Wort der Warnung von mir: Achtet stets auf das charakteristische Schimmern! 

Einer alten Legende nach thronen die Flussjungfern in einsamen Nächten auf bemoosten Felsen und richten ihr endloses, wasserfarbenes Haar mit Kämmen aus Knochen und altem Treibholz. Ihr Antlitz sei dabei von einer friedlichen Melancholie gezeichnet, während sie leise Lieder summen, deren tiefere Bedeutung allein dem fließenden Wasser vorbehalten bleibt.

    \begin{tabular}{@{}ll@{}}
        \textbf{\trans{weight_label}:} & 0.10 \\
        \textbf{\trans{price_label}:} & 25 \\
        \textbf{\trans{rarity_label}:} & Außergewöhnlich \\
        \textbf{\trans{concealment_label}:} & 0 \\
    \end{tabular}

    \wbif{phasesix}{
        \vspace{3mm}
\faIcon{flask}\hspace{0.5em}    }{}


    \subsection*{Silberorchidee}

    Die Silberorchidee gilt als die unangefochtene und trügerische Königin der südlichen Flora. Sie ist ein botanisches Wunderwerk, das ebenso schön wie tödlich für jene ist, die sich von ihrem Glanz täuschen lassen.

Ihre Blätter sind nicht grün, sondern besitzen eine dunkelgraue, fast metallische Färbung, die im Mondlicht wie poliertes Silber glänzt. Die Adern pulsieren schwach in einem blassen Violett, wenn Magie in der Nähe ist. Die Blüte selbst ist groß und kelchförmig, mit schneeweißen Blütenblättern, die am Rand rasiermesserscharf sind. Doch das Verstörendste an ihr ist nicht ihre Schönheit, sondern ihre Mobilität: Die Pflanze reckt sich auf freiliegenden, muskulösen Wurzeln empor, mit denen sie in der Lage ist, langsam kriechend über den Boden zu wandern.

Die Silberorchidee ist fast ausschließlich im tiefen Süden Tirakans heimisch, jenseits der Eisernen Berge. Man findet sie in den weiten, grünen Steppen und an Flussläufen, oft im Schatten der dortigen Flora getarnt. Sie wächst oft in beunruhigender Nähe zu den riesenhaften Wesen aus Stein, die in den Pässen hausen.

Nähert man sich der Pflanze, stößt sie eine glitzernde Wolke aus feinem, silbernem Staub aus. Dies ist kein harmloser Pollenflug, sondern ein tödlicher Angriff. Wer den Staub einatmet, wird von schwerem Husten und Atemnot ergriffen. Binnen Augenblicken bilden sich schwarze Pocken auf der Haut, und das Opfer fällt in eine tiefe, todesähnliche Ohnmacht. Ist das Opfer wehrlos, sondert die Pflanze ein ätzendes Sekret aus ihren Wurzeln ab, um die Beute langsam zu zersetzen und aufzunehmen.

Trotz dieser Gefahren wird sie gejagt, denn sie ist ein mächtiger Potenzierer. In der Alchemie wird der extrahierte Nektar genutzt, um die Wirkung anderer Tränke ins Extreme zu steigern. Doch die Verarbeitung ist riskant: Ein Fehler in der Destillation führt dazu, dass die magische Energie den Körper überlädt (analog zum Silbertod), was oft tödlich endet.

In den alten Legenden heißt es, die Silberorchideen seien entstanden, als im Ersten Zeitalter das Blut eines gefallenen Sternengottes auf die Erde tropfte. Die Elfen hingegen nennen die Blume Verräterschmuck und glauben, dass sie dort wächst, wo die Realität Risse bekommen hat und Chaos in die Welt sickert.

Sieht aus wie Schmuck, kriecht wie eine Spinne und ist mehr wert als mein Haus. Aber pass auf, Junge: Wenn du das Glitzern siehst, halte den Atem an und renn. Bevor du merkst, dass du krank wirst, bist du schon ihr Dünger. – Randnotiz in den Aufzeichnungen des Kräutersammlers 'Drei-Finger-Hannes'

    \begin{tabular}{@{}ll@{}}
        \textbf{\trans{weight_label}:} & 0.10 \\
        \textbf{\trans{price_label}:} & 550 \\
        \textbf{\trans{rarity_label}:} & Selten \\
        \textbf{\trans{concealment_label}:} & 0 \\
    \end{tabular}

    \wbif{phasesix}{
        \vspace{3mm}
\faIcon{flask}\hspace{0.5em}    }{}


    \subsection*{Lavendel (Lavandula officinalis)}

    Im elften Jahrhundert wurde Lavendel von Mönchen in Mitteleuropa angesiedelt. In der Medizin wurde Lavendel eine Wirksamkeit bei Insektenstichen und Verbrennungen zugeschrieben. Ein Lavendel Tee hilft bei Erkältungen und Kopfschmerzen.

    \begin{tabular}{@{}ll@{}}
        \textbf{\trans{weight_label}:} & 0.10 \\
        \textbf{\trans{price_label}:} & 4 \\
        \textbf{\trans{rarity_label}:} & Gewöhnlich \\
        \textbf{\trans{concealment_label}:} & 0 \\
    \end{tabular}

    \wbif{phasesix}{
        \vspace{3mm}
\faIcon{cube}\hspace{0.5em}\faIcon{flask}\hspace{0.5em}    }{}


    \subsection*{Totenmohn}

    Der Totenmohn wird ausschließlich auf der asgoraner Insel Linya angebaut, das Gewächs lässt sich sonst nirgends kultivieren. Bei der Pflanze handelt es sich um eine etwa zwanzig Finger hohe, mohnartige Blüte, welche zum Teil schwarz gefärbt ist. 

Der Totenmohn wird für allerhand Rituale und Tränke verwendet, was die Pflanze zu einem wahren Exportschlager Asgorans werden lässt.

    \begin{tabular}{@{}ll@{}}
        \textbf{\trans{weight_label}:} & 0.10 \\
        \textbf{\trans{price_label}:} & 20 \\
        \textbf{\trans{rarity_label}:} & Außergewöhnlich \\
        \textbf{\trans{concealment_label}:} & 0 \\
    \end{tabular}

    \wbif{phasesix}{
        \vspace{3mm}
\faIcon{dragon}\hspace{0.5em}\faIcon{flask}\hspace{0.5em}    }{}


    \subsection*{Sonnenblüte}

    Die Sonnenblüte ist der Inbegriff von Beständigkeit. Mit ihrem kräftigen, rauen Stängel und dem stolzen, goldgelben Kranz aus Blütenblättern folgt sie unermüdlich dem Lauf des Sonnenwagens über das Firmament von Tiraka.

Ihr Kern ist gefüllt mit nahrhaften, ölhaltigen Samen, doch für Alchemisten sind es vor allem die leuchtenden Randblätter, die den Wert ausmachen. Sie speichern die reine, sanfte Wärme des Tages, ohne dabei die gefährliche Hitze des Feuers oder die Unberechenbarkeit der Magie in sich zu tragen.

Man findet sie fast überall in den fruchtbaren Ebenen Tirakas, besonders auf den sonnenverwöhnten Hügeln rund um Asgoran oder in den Gärten der Bauern im Hinterland. Sie liebt offene Flächen und tiefen, schwarzen Boden. Es ist keine seltene Pflanze, nein, aber eine, die Pflege braucht. Die wilden Varianten in den Heiden sind oft kleiner, aber ihre Essenz ist konzentrierter als die der gezüchteten Prachtexemplare.

Die Blätter sollten zur Mittagsstunde gezupft werden, wenn die Sonne am höchsten steht und die Blüte ihre volle Pracht entfaltet hat. Man trocknet sie flach liegend auf Leinentüchern an einem luftigen Ort, niemals im direkten Ofenfeuer!

    \begin{tabular}{@{}ll@{}}
        \textbf{\trans{weight_label}:} & 0.10 \\
        \textbf{\trans{price_label}:} & 2 \\
        \textbf{\trans{rarity_label}:} & Gewöhnlich \\
        \textbf{\trans{concealment_label}:} & 0 \\
    \end{tabular}

    \wbif{phasesix}{
        \vspace{3mm}
\faIcon{flask}\hspace{0.5em}    }{}


    \subsection*{Tran des Schillerwals}

    Der Tran, das Fettgewebe unter der Haut, ist der Grund, warum diese majestätischen Tiere gejagt werden. Es ist kein gewöhnliches Fett, sondern ein  Speicher magischer Energie.

Der Tran ist im rohen Zustand eine zähe, geleeartige Masse, die schwach bläulich leuchtet. Nach der Veredelung (Schmelzen und Filtern) wird er zu einem klaren, öligen Elixier, das wie flüssiges Perlmutt schlieren zieht.
Im Gegensatz zum ranzigen Tran gewöhnlicher Wale riecht der Tran des Schillerwals frisch, salzig und leicht metallisch (wie die Luft vor einem Gewitter).

Er ist das beste bekannte Mittel, um flüchtige Magie in Tränken zu binden (siehe Elixier der Elfenwacht).

In Lampen verbrannt, gibt er ein Licht ab, das niemals rußt und unsichtbares sichtbar machen kann. Waffenöle aus diesem Tran lassen Geister verletzen.

    \begin{tabular}{@{}ll@{}}
        \textbf{\trans{weight_label}:} & 0.10 \\
        \textbf{\trans{price_label}:} & 600 \\
        \textbf{\trans{rarity_label}:} & Selten \\
        \textbf{\trans{concealment_label}:} & 0 \\
    \end{tabular}

    \wbif{phasesix}{
        \vspace{3mm}
\faIcon{dragon}\hspace{0.5em}\faIcon{flask}\hspace{0.5em}    }{}


    \subsection*{Arnika (Arnica montana)}

    Arnika wird bei Entzündungen, Wunden, um den Kreislauf anzuregen und als Abtreibungsmittel verwendet. Die Blüten werden als Salbe, als Tee oder als Tinktur verwendet.

    \begin{tabular}{@{}ll@{}}
        \textbf{\trans{weight_label}:} & 0.10 \\
        \textbf{\trans{price_label}:} & 5 \\
        \textbf{\trans{rarity_label}:} & Gewöhnlich \\
        \textbf{\trans{concealment_label}:} & 0 \\
    \end{tabular}

    \wbif{phasesix}{
        \vspace{3mm}
\faIcon{cube}\hspace{0.5em}    }{}


    \subsection*{Wasser}

    Kaltes klares Wasser.

    \begin{tabular}{@{}ll@{}}
        \textbf{\trans{weight_label}:} & 1.00 \\
        \textbf{\trans{price_label}:} & 0 \\
        \textbf{\trans{rarity_label}:} & Gewöhnlich \\
        \textbf{\trans{concealment_label}:} & 0 \\
    \end{tabular}

    \wbif{phasesix}{
        \vspace{3mm}
\faIcon{flask}\hspace{0.5em}    }{}


    \subsection*{Alkohol}

    In den Laboren der Alchemisten und den Hütten der Kräuterkundigen wird Alkohol selten zum Vergnügen gelagert. Er gilt als Lösungsmittel, dazu fähig, die Essenz von Pflanzen und Mineralien zu entfesseln.

Er wird durch mehrfache Destillation aus vergorenem Getreide oder Früchten gewonnen. Esist ein brennend scharf riechendes Destillat, das so stark ist, dass es beim Einatmen die Nase kitzelt und auf der Haut sofort kühl verdunstet.

    \begin{tabular}{@{}ll@{}}
        \textbf{\trans{weight_label}:} & 0.50 \\
        \textbf{\trans{price_label}:} & 1 \\
        \textbf{\trans{rarity_label}:} & Gewöhnlich \\
        \textbf{\trans{concealment_label}:} & 0 \\
    \end{tabular}

    \wbif{phasesix}{
        \vspace{3mm}
\faIcon{flask}\hspace{0.5em}    }{}


    \subsection*{Felsentroll-Klauen (Lithocrinus tirakanis)}

    Die Klauen eines Felsentrolls sind keine Krallen im biologischen Sinne, sondern mineralisierte Auswüchse aus verhärtetem Keratin und konzentrierten Erzen.

Sie spiegeln die unbändige, physische Regenerationskraft der Trolle wider, deren Haut eine Symbiose mit dem Gestein der Schattenfelsen eingegangen ist.

Eine einzelne, intakte Klaue hat etwa die Größe eines Kurzschwerters.
Die Färbung reicht von tiefem Grau bis hin zu Obsidian-Schwarz, wobei oft Schichten erkennbar sind, die an Schiefer erinnern.

Sie sind so hart, dass sie Funken schlagen, wenn sie auf Metall treffen. Gewöhnliche Klingen stumpfen an ihnen meistens sofort ab.

In der Alchemie wird fast nie die ganze Klaue verwendet, sondern eine verarbeitete Form. Die Klaue muss mühsam mit Diamantfeilen oder Meißeln aus gehärtetem Stahl bearbeitet werden. 

Als Zutat dienen fein geraspelte Späne oder Staub. Dieser Staub ist schwer, aschefarben und glitzert unter Lichteinfall.
Der Staub muss extrem fein sein. Zu grobe Späne lösen sich im Trank nicht auf und können beim Konsum die Speiseröhre des Anwenders schwer verletzen. Alternativ ist es ratsam, nach dem Brauvorgang, die Flüssigkeit durch ein Sieb zu geben.

    \begin{tabular}{@{}ll@{}}
        \textbf{\trans{weight_label}:} & 1.00 \\
        \textbf{\trans{price_label}:} & 125 \\
        \textbf{\trans{rarity_label}:} & Selten \\
        \textbf{\trans{concealment_label}:} & 0 \\
    \end{tabular}

    \wbif{phasesix}{
        \vspace{3mm}
\faIcon{flask}\hspace{0.5em}    }{}


    \subsection*{Silberspähne}

    In Tiraka werden Silberspäne meist als Nebenprodukt in Schmieden oder bei der Herstellung von Schmuck gewonnen.

Für alchemistische Zwecke werden sie oft im Feuer gereinigt, um als Läutersilber die magischen Ströme in Tränken zu stabilisieren.
Silberspäne dienen als energetischer Anker. Sie verhindern, dass die instabilen Komponenten den Trank während des Brauprozesses zerreißen.

    \begin{tabular}{@{}ll@{}}
        \textbf{\trans{weight_label}:} & 0.10 \\
        \textbf{\trans{price_label}:} & 1 \\
        \textbf{\trans{rarity_label}:} & Gewöhnlich \\
        \textbf{\trans{concealment_label}:} & 0 \\
    \end{tabular}

    \wbif{phasesix}{
        \vspace{3mm}
\faIcon{flask}\hspace{0.5em}    }{}


    \subsection*{Goldnugget}

    Ein kleines Stück unverarbeitetes Gold, etwa 5 Gramm.

    \begin{tabular}{@{}ll@{}}
        \textbf{\trans{weight_label}:} & 0.05 \\
        \textbf{\trans{price_label}:} & 300 \\
        \textbf{\trans{rarity_label}:} & Gewöhnlich \\
        \textbf{\trans{concealment_label}:} & 0 \\
    \end{tabular}

    \wbif{phasesix}{
        \vspace{3mm}
\faIcon{cube}\hspace{0.5em}    }{}


    \subsection*{Phiole Harpyenblut}

    Eine Phiole gefüllt mit Harpyenblut.

    \begin{tabular}{@{}ll@{}}
        \textbf{\trans{weight_label}:} & 0.10 \\
        \textbf{\trans{price_label}:} & 5 \\
        \textbf{\trans{rarity_label}:} & Selten \\
        \textbf{\trans{concealment_label}:} & 0 \\
    \end{tabular}

    \wbif{phasesix}{
        \vspace{3mm}
\faIcon{flask}\hspace{0.5em}\faIcon{magic}\hspace{0.5em}    }{}


    \subsection*{Trollzahn}

    Der Fangzahn eines Trolls

    \begin{tabular}{@{}ll@{}}
        \textbf{\trans{weight_label}:} & 0.10 \\
        \textbf{\trans{price_label}:} & 30 \\
        \textbf{\trans{rarity_label}:} & Selten \\
        \textbf{\trans{concealment_label}:} & 0 \\
    \end{tabular}

    \wbif{phasesix}{
        \vspace{3mm}
\faIcon{flask}\hspace{0.5em}\faIcon{magic}\hspace{0.5em}    }{}


    \subsection*{Nachtschiefer-Staub}

    Der Nachtschiefer ist kein gewöhnliches Mineral, sondern die fein zermahlene Essenz uralter, tiefschwarzer Gesteinsadern.

Dieses Material wird nicht geschürft, sondern muss aus dem Herz von Bergwerken gewonnen werden, die seit Äonen kein Sonnenlicht mehr gesehen haben, oft nur zugänglich durch enge Stollen in den kältesten, abgelegensten Gebirgszügen.

Man findet ihn ausschließlich in tiefen, stillgelegten Minen oder unterirdischen Krypten, wo das Gestein über Jahrtausende hinweg durch den konstanten Druck und die absolute Kälte des Erdinneren verdichtet wurde. Die Adern schimmern leicht, wenn sie im Schein einer Fackel erfasst werden. Ein Zeichen ihrer beinahe unnatürlichen Reinheit.

    \begin{tabular}{@{}ll@{}}
        \textbf{\trans{weight_label}:} & 1.00 \\
        \textbf{\trans{price_label}:} & 40 \\
        \textbf{\trans{rarity_label}:} & Selten \\
        \textbf{\trans{concealment_label}:} & 0 \\
    \end{tabular}

    \wbif{phasesix}{
        \vspace{3mm}
\faIcon{flask}\hspace{0.5em}    }{}


    \subsection*{Melisse (Melissa officinalis)}

    Melisse wurde seit jeher als Heilkraut in der Medizin genutzt. Sie wirkt gegen Kopfschmerzen, Nervosität, Schlafstörungen und Magen-Darm-Beschwerden. Zudem bringt ein Aufguss mit Melisse Entspannung.

    \begin{tabular}{@{}ll@{}}
        \textbf{\trans{weight_label}:} & 0.10 \\
        \textbf{\trans{price_label}:} & 2 \\
        \textbf{\trans{rarity_label}:} & Gewöhnlich \\
        \textbf{\trans{concealment_label}:} & 0 \\
    \end{tabular}

    \wbif{phasesix}{
        \vspace{3mm}
\faIcon{cube}\hspace{0.5em}\faIcon{flask}\hspace{0.5em}    }{}


    \subsection*{Engelwurzen (Angelica archangelica)}

    Die Pflanze wird bei Verdauungsschwäche, Appetitlosigkeit und Verdauungsschwäche angewendet, und soll angeblich vor der Pest schützen.

    \begin{tabular}{@{}ll@{}}
        \textbf{\trans{weight_label}:} & 0.10 \\
        \textbf{\trans{price_label}:} & 3 \\
        \textbf{\trans{rarity_label}:} & Gewöhnlich \\
        \textbf{\trans{concealment_label}:} & 0 \\
    \end{tabular}

    \wbif{phasesix}{
        \vspace{3mm}
\faIcon{cube}\hspace{0.5em}\faIcon{flask}\hspace{0.5em}    }{}


    \subsection*{Gildenkraut (Vigilia mercatoris)}

    Das Gildenkraut ist eine zähe, beinahe unscheinbare Pflanze, die nur dem geschulten Auge eines Kundigen auffällt.

Sie zeichnet sich durch gezackte, ledrige Blätter von tiefgrüner Färbung aus, deren dicke Blattadern einen leichten rötlichen Schimmer aufweisen. Ihre Blüten sind winzig, blassgelb und schließen sich sofort, sobald die Sonne ihren Zenit überschreitet.

Das Kraut ist vollkommen mundan und besitzt nicht den geringsten Funken Arkana. Sein Wert liegt vielmehr in der extrem hohen, natürlichen Konzentration an Alkaloiden, vor allem reinem Koffein. Zerreibt man die harten Blätter zwischen den Fingern, verströmen sie einen scharfen, erdigen Duft, der an stark geröstete Nüsse erinnert.

Lage und Lebensraum:
Dieses Gewächs ist verhältnismäßig selten, da es ganz bestimmte, widrige Umstände benötigt, um seine anregenden Säfte zu bilden. Es gedeiht ausschließlich in kargen, mineralreichen Böden und meidet das feuchte Moos tiefer Wälder.

Man findet das Gildenkraut vornehmlich in den felsigen Spalten uralter Handelsrouten, im bröckelnden Steinmörtel verlassener Gilden-Außenposten oder an den steilen, sonnengegerbten Hängen des Toraner Hochlands. Es braucht den harten Fels und viel Wind, um nicht zu verkümmern.

    \begin{tabular}{@{}ll@{}}
        \textbf{\trans{weight_label}:} & 1.00 \\
        \textbf{\trans{price_label}:} & 50 \\
        \textbf{\trans{rarity_label}:} & Selten \\
        \textbf{\trans{concealment_label}:} & 0 \\
    \end{tabular}

    \wbif{phasesix}{
        \vspace{3mm}
\faIcon{dragon}\hspace{0.5em}\faIcon{flask}\hspace{0.5em}    }{}


    \subsection*{Phiole Flugechsenblut}

    Meistens extrahiert aus den Adern der Flugechsen, die von den O'Gru domestiziert wurden. Es unterscheidet sich nicht vom Blut der in der Wildnis lebenden Exemplare.

Eine Phiole enthält 50 Oth Flugechselnblut.

    \begin{tabular}{@{}ll@{}}
        \textbf{\trans{weight_label}:} & 0.10 \\
        \textbf{\trans{price_label}:} & 30 \\
        \textbf{\trans{rarity_label}:} & Außergewöhnlich \\
        \textbf{\trans{concealment_label}:} & 0 \\
    \end{tabular}

    \wbif{phasesix}{
        \vspace{3mm}
\faIcon{flask}\hspace{0.5em}    }{}


    \subsection*{Alant (Inula helenium)}

    Diese Heilpflanze aus dem Mittelalter ist in modernen Zeiten nicht mehr weit verbreitet. Ihre Anwendung verbessert die Verdauung, und ihr wird eine vorbeugende Wirkung gegen Darmkrebs zugeschrieben.

    \begin{tabular}{@{}ll@{}}
        \textbf{\trans{weight_label}:} & 0.10 \\
        \textbf{\trans{price_label}:} & 5 \\
        \textbf{\trans{rarity_label}:} & Gewöhnlich \\
        \textbf{\trans{concealment_label}:} & 0 \\
    \end{tabular}

    \wbif{phasesix}{
        \vspace{3mm}
\faIcon{cube}\hspace{0.5em}\faIcon{flask}\hspace{0.5em}    }{}


    \subsection*{Bergwyvern Galle}

    Die Galle ist eine hochviskose, leuchtend goldgelbe bis giftgrüne Flüssigkeit, die in der Gallenblase der Bergwyvern gespeichert wird.

Der Geruch ist beißend, schwefelig und mit einer Note von verbranntem Kupfer. Schon das Einatmen der reinen Dämpfe kann die Nasenschleimhäute verätzen.

In Tiraka glaubt man, dass die Galle die „Essenz des ungestillten Hungers“ enthält. In der Alchemie wird sie als Katalysator verwendet, um andere Zutaten gewaltsam miteinander zu verschmelzen, die sich normalerweise abstoßen würden.

Die Gewinnung der Galle ist ein schwieriges Unterfangen für einen Alchemisten, da sie höchste Präzision unter widrigen Bedingungen erfordert.

Die Galle muss innerhalb von 1W6 Stunden nach dem Tod der Kreatur entnommen werden. Danach beginnt die Zersetzung der Gallenblase, und die Flüssigkeit verliert ihre alchemistische Potenz.

Man benötigt chirurgisches Besteck aus gehärtetem Stahl oder speziellen Keramikmessern. Einfaches Eisen würde von der Säure binnen Sekunden zerfressen werden.

Der Prozess des Eingriffs läuft wie folgt ab: Der Kadaver muss auf dem Rücken fixiert werden.
Ein tiefer Schnitt unterhalb des Brustbeins legt die Leber frei.
Die Gallenblase ist ein prall gefüllter, pulsierender Beutel. Sie muss mit einer Klemme am oberen Ende abgebunden werden, bevor man sie vorsichtig herausschneidet. Ein erfolgreicher Wurf auf Geschick (oder Medizin) gegen MW 8 ist erforderlich.

Sollte der Beutel platzen, erleidet der Erntende sofort 1W6 Schaden durch Verätzung, und die Zutat ist unwiederbringlich verloren. 

Wyvern-Galle kann nicht in normalen Glasfläschchen aufbewahrt werden, da sie das Glas mit der Zeit „blind“ macht und spröde werden lässt. Erfahrene Abenteurer nutzen bleigefütterte Tonkrüge oder Gefäße aus reinem Quarz, um die Substanz sicher nach Hause zu bringen.

    \begin{tabular}{@{}ll@{}}
        \textbf{\trans{weight_label}:} & 2.50 \\
        \textbf{\trans{price_label}:} & 500 \\
        \textbf{\trans{rarity_label}:} & Selten \\
        \textbf{\trans{concealment_label}:} & 1 \\
    \end{tabular}

    \wbif{phasesix}{
        \vspace{3mm}
\faIcon{flask}\hspace{0.5em}    }{}


    \subsection*{Beinwell (Symphytum officinale)}

    Beinwell regt die Durchblutung an, Prellungen, Hämatome und Verstauchungen verschwinden schneller. Beinwell beschleunigt die Regeneration der Zellen.

    \begin{tabular}{@{}ll@{}}
        \textbf{\trans{weight_label}:} & 0.10 \\
        \textbf{\trans{price_label}:} & 5 \\
        \textbf{\trans{rarity_label}:} & Gewöhnlich \\
        \textbf{\trans{concealment_label}:} & 0 \\
    \end{tabular}

    \wbif{phasesix}{
        \vspace{3mm}
\faIcon{cube}\hspace{0.5em}    }{}


    \subsection*{Thymian (Thymus vulgaris)}

    Thymian wird seit über 4000 Jahren gegen Keuchhusten, Husten und Bronchitis eingesetzt. Seine schleimlösende Wirkung wird besonders geschätzt.

    \begin{tabular}{@{}ll@{}}
        \textbf{\trans{weight_label}:} & 0.10 \\
        \textbf{\trans{price_label}:} & 5 \\
        \textbf{\trans{rarity_label}:} & Gewöhnlich \\
        \textbf{\trans{concealment_label}:} & 0 \\
    \end{tabular}

    \wbif{phasesix}{
        \vspace{3mm}
\faIcon{cube}\hspace{0.5em}\faIcon{flask}\hspace{0.5em}    }{}


    \subsection*{Schöllkraut (Chelidonium majus)}

    Im Mittelalter wurde Schöllkraut bei Hautausschlägen, Sehschwäche oder Gelbsucht verwendet. Die Alkaloide der Pflanze habenden einen krampflösenden Effekt. Sie helfen bei Verdauungsproblemen und regen den Gallenfluss an.

    \begin{tabular}{@{}ll@{}}
        \textbf{\trans{weight_label}:} & 0.10 \\
        \textbf{\trans{price_label}:} & 5 \\
        \textbf{\trans{rarity_label}:} & Gewöhnlich \\
        \textbf{\trans{concealment_label}:} & 0 \\
    \end{tabular}

    \wbif{phasesix}{
        \vspace{3mm}
\faIcon{cube}\hspace{0.5em}\faIcon{flask}\hspace{0.5em}    }{}


    \subsection*{Kinstarchel Sekret}

    Eine hoch reaktive Flüssigkeit, die aus den Knochen der possierlichen Kinstarchel hergestellt wird.

    \begin{tabular}{@{}ll@{}}
        \textbf{\trans{weight_label}:} & 0.10 \\
        \textbf{\trans{price_label}:} & 30 \\
        \textbf{\trans{rarity_label}:} & Selten \\
        \textbf{\trans{concealment_label}:} & 0 \\
    \end{tabular}

    \wbif{phasesix}{
        \vspace{3mm}
\faIcon{dragon}\hspace{0.5em}\faIcon{flask}\hspace{0.5em}    }{}


    \subsection*{Kiefernharz}

    Das schützende Blut des Baumes. Wenn die raue Rinde einer Kiefer,sei es durch den Biss eines Wildtieres oder das ungeschickte Beil eines Holzfällers, verletzt wird, tritt diese zähe, goldgelbe Flüssigkeit aus.

Sie fließt langsam, beinahe träge, und erfüllt die Umgebung mit diesem unverwechselbaren, herben Duft.

Es wird genutzt, um spröde Korken zu versiegeln, Fackeln eine kräftige Flamme zu verleihen oder um einfache Wundsalben dickflüssig zu machen.

Sobald es an der Luft trocknet, wird es steinhart und schließt die Wunde des Baumes wie ein natürlicher Pfropfen. 

Wenn Ihr das Zeug erst einmal an den Händen habt, hilft nur noch ordentlich Fett oder sehr viel Geduld, um den Kleber wieder loszuwerden.

    \begin{tabular}{@{}ll@{}}
        \textbf{\trans{weight_label}:} & 0.50 \\
        \textbf{\trans{price_label}:} & 2 \\
        \textbf{\trans{rarity_label}:} & Gewöhnlich \\
        \textbf{\trans{concealment_label}:} & 0 \\
    \end{tabular}

    \wbif{phasesix}{
        \vspace{3mm}
\faIcon{flask}\hspace{0.5em}    }{}


    \subsection*{Baldrian (Valeriana officinalis)}

    Baldrian hilft bei Einschlafstörungen und Unruhe. Hopfen und Melisse steigern die Wirkung des Baldrians und verbessern den Geschmack.

    \begin{tabular}{@{}ll@{}}
        \textbf{\trans{weight_label}:} & 0.10 \\
        \textbf{\trans{price_label}:} & 3 \\
        \textbf{\trans{rarity_label}:} & Gewöhnlich \\
        \textbf{\trans{concealment_label}:} & 0 \\
    \end{tabular}

    \wbif{phasesix}{
        \vspace{3mm}
\faIcon{cube}\hspace{0.5em}\faIcon{flask}\hspace{0.5em}    }{}


    \subsection*{Kinstarchel Knochen}

    Knochen eines verstorbenen Kinstarchels. Aus dem Mark der Knochen kann ein explosives Sekret gewonnen werden.

    \begin{tabular}{@{}ll@{}}
        \textbf{\trans{weight_label}:} & 1.00 \\
        \textbf{\trans{price_label}:} & 25 \\
        \textbf{\trans{rarity_label}:} & Außergewöhnlich \\
        \textbf{\trans{concealment_label}:} & 0 \\
    \end{tabular}

    \wbif{phasesix}{
        \vspace{3mm}
\faIcon{flask}\hspace{0.5em}    }{}


    \subsection*{Taniumstaub}

    Tanium ist ein dunkles, kristallines Element von außergewöhnlicher Härte. Im Rohzustand findet man es oft als tiefschwarze Adern in uraltem Gestein, meist in Gebieten mit hoher natürlicher Magiekonzentration.

Es ist so hart, dass herkömmliche Mörser daran zerbrechen; nur Werkzeuge aus gehärtetem Diamant oder magisch verstärkte Mahlwerke können es zu feinem Staub zermahlen.

Tanium fungiert als perfekter Speicher für arkane Energie. Doch Vorsicht: Es hat kein Sättigungsgefühl! Wird es mit zu viel Magie gesättigt oder durch unsaubere Alchemie instabil, entlädt es sich in einer magischen Explosion.

    \begin{tabular}{@{}ll@{}}
        \textbf{\trans{weight_label}:} & 0.10 \\
        \textbf{\trans{price_label}:} & 40 \\
        \textbf{\trans{rarity_label}:} & Außergewöhnlich \\
        \textbf{\trans{concealment_label}:} & 0 \\
    \end{tabular}

    \wbif{phasesix}{
        \vspace{3mm}
\faIcon{flask}\hspace{0.5em}    }{}


    \subsection*{Spitzwegerich (Plantago lanceolata)}

    Die spitzen, schmalen Blätter des Spitzwegerich werden als Sirup oder auch als Tee bei Erkältungskrankheiten verwendet. Spitzwegerich kann auch zerkleinert und zerrieben auf Wunden oder Insektenstiche aufgetragen werden und entfaltet hier eine kühlende Wirkung. Die Pflanze wird auch bei Durchfall eingesetzt.

    \begin{tabular}{@{}ll@{}}
        \textbf{\trans{weight_label}:} & 0.10 \\
        \textbf{\trans{price_label}:} & 2 \\
        \textbf{\trans{rarity_label}:} & Gewöhnlich \\
        \textbf{\trans{concealment_label}:} & 0 \\
    \end{tabular}

    \wbif{phasesix}{
        \vspace{3mm}
\faIcon{cube}\hspace{0.5em}\faIcon{flask}\hspace{0.5em}    }{}


    \section{Musikinstrumente}



    \subsection*{Laute}

    Die Laute (über spätmittelhochdeutsch lūte von arabisch العود, DMG al-ʿūd ‚der Stab, das Holz, Laute‘) ist ein Zupfinstrument mit Korpus und angesetztem Hals sowie mit gleichlaufend zur Instrumentendecke verlaufenden Saiten.

    \begin{tabular}{@{}ll@{}}
        \textbf{\trans{weight_label}:} & 1.50 \\
        \textbf{\trans{price_label}:} & 25 \\
        \textbf{\trans{rarity_label}:} & Gewöhnlich \\
        \textbf{\trans{concealment_label}:} & 3 \\
    \end{tabular}

    \wbif{phasesix}{
        \vspace{3mm}
\faIcon{eye}\hspace{0.5em}\faIcon{chess-rook}\hspace{0.5em}\faIcon{umbrella}\hspace{0.5em}\faIcon{flag}\hspace{0.5em}\faIcon{plane}\hspace{0.5em}\faIcon{microchip}\hspace{0.5em}\faIcon{ankh}\hspace{0.5em}\faIcon{9}\hspace{0.5em}    }{}


    \subsection*{Dudelsack}

    Zwergischer Dudelsack. Da ist genug Luft im Behälter um zu Singen während Zwerg spielt.

    \begin{tabular}{@{}ll@{}}
        \textbf{\trans{weight_label}:} & 5.00 \\
        \textbf{\trans{price_label}:} & 250 \\
        \textbf{\trans{rarity_label}:} & Gewöhnlich \\
        \textbf{\trans{concealment_label}:} & 0 \\
    \end{tabular}

    \wbif{phasesix}{
        \vspace{3mm}
\faIcon{chess-rook}\hspace{0.5em}\faIcon{umbrella}\hspace{0.5em}\faIcon{flag}\hspace{0.5em}\faIcon{plane}\hspace{0.5em}\faIcon{microchip}\hspace{0.5em}\faIcon{dragon}\hspace{0.5em}\faIcon{magic}\hspace{0.5em}\faIcon{spider}\hspace{0.5em}\faIcon{fire}\hspace{0.5em}    }{}


    \subsection*{Flöte der Erleuchtung}

    Eine kunstvoll gearbeitete Holzflöte, verziert mit filigranen Mustern, die von ihrer besonderen Bedeutung zeugen. Damit sie ihren Zauber entfalten kann, ist eine gelungene Darbietungsprobe erforderlich. Beim Spielen inspiriert sie den Spieler, erfüllt ihn mit schöpferischer Energie und verringert den Mindestwurf für zwei Runden um 1.

    \begin{tabular}{@{}ll@{}}
        \textbf{\trans{weight_label}:} & 0.20 \\
        \textbf{\trans{price_label}:} & 150 \\
        \textbf{\trans{rarity_label}:} & Einmalig \\
        \textbf{\trans{concealment_label}:} & 0 \\
    \end{tabular}

    \wbif{phasesix}{
        \vspace{3mm}
\faIcon{dragon}\hspace{0.5em}    }{}

\end{multicols}